\documentclass[Total.tex]{subfiles}

\begin{document}
	Examples Reference:\url{https://adjectivesproject.org/?page=data}
	\chapter{Affine Schemes}
		\section{Locally Ringed Spaces}
		
			\textbf{Note} The knowledge here should be listed in the part of sheaves. But the locally ringed is mainly used in algebraic geometry. It is convenient to introduce here. 
			
			\begin{Def}
				\begin{itemize}
					\item A \textbf{locally ringed space} is sheaf $(X,\mathcal{O}_X)$, whose stalks $\mathcal{O}_{X,x},x\in X$ are local rings.
					\item Define the \textbf{residue field}
					\begin{gather*}
						\kappa(x):=\mathcal{O}_{X,x}/m_x
					\end{gather*}
					\item The morphism of locally ringed spaces is the morphism of ringed spaces
					\begin{gather*}
						(f,f^\#):(X,\mathcal{O}_X)\rightarrow (Y,\mathcal{O}_Y)
					\end{gather*} 
					via this lemma: \ref{Sheaves-Sheaves-Equivalence-of-Morphsims-Induced-by-Continuous-Maps} There exists a correspondence(It is actually an $f$-map)
					\begin{gather*}
						f^\#:f^{-1}\mathcal{O}_Y\rightarrow \mathcal{O}_X\qquad f^\flat:\mathcal{O}_Y\rightarrow f_*\mathcal{O}_X
					\end{gather*}
					\item Composition: Consider 
					\begin{gather*}
						f:X\rightarrow Y\qquad g:Y\rightarrow Z
					\end{gather*}
					
					Then, there exists a composition
					\begin{gather*}
						\mathcal{O}_{Z,g(f(x))}\stackrel{g^\#}{\rightarrow} \mathcal{O}_{Y,f(x)}\stackrel{f^\#}{\rightarrow} \mathcal{O}_{X,x}
					\end{gather*}
				\end{itemize}
			\end{Def}
			
			Then, $(f,f^\#)$ induces the morphism of stalks
			\begin{gather*}
				f^\#_x:\mathcal{O}_{Y,f(x)}\rightarrow \mathcal{O}_{X,X}
			\end{gather*}
			
			which is a local ring map. It maps $m_{Y,f(x)}$ to $m_{X,x}$.
			
			It is easy to check, stalks isomorphism $\Leftrightarrow$ Local ring isomorphisms.
			
			
			\begin{Def}
				Trivial: Functor
				\begin{center}
					Locally ringed spaces$\rightarrow$ Ringed Spaces
				\end{center}
				
				which reflects isomorphisms. 
			\end{Def}
			\subsection{Morphisms}
			\subsubsection{Open Immersion of Locally Ringed Spaces}
				
				\begin{Def}
					Let $f:X\rightarrow Y$ be a morphism of locally ringed spaces. We say that $f$ is an \textbf{open immersion} if $f$ is  homeomorphism of $X$ onto an open subset of $Y$, and the map
					\begin{gather*}
						f^\#:f^{-1}\mathcal{O}_Y\rightarrow \mathcal{O}_X
					\end{gather*}
					is an isomorphism.
				\end{Def}
				\begin{example}
					Consider the open subset $U\subseteq X$, then 
					\begin{gather*}
						j:(U,\mathcal{O}_U)\rightarrow (X,\mathcal{O}_X)
					\end{gather*}
					
					is an open immersion. Where $\mathcal{O}_U(V):=\mathcal{O}_X(V);V\subseteq U$.
				\end{example}
				\begin{Def}
					Let $X$ be a locally ringed space.$U\subseteq X$ is a closed subset. Then, the scheme above is called an \textbf{open subspace} of scheme $(X,\mathcal{O}_X)$.
				\end{Def}
				\begin{lemma}
					Let $f:X\rightarrow Y$ be an open immersion of locally ringed spaces with $j:V=f(X)\rightarrow Y$ be the open subspace of $Y$ assocaited to the image of $f$. Then: There exists unique isomorphism $f':X\cong V$ of locally ringed spaces such that
					\begin{gather*}
						f=j\circ f'
					\end{gather*}
				\end{lemma}
				\begin{proof}
					Firstly, let $f':X\rightarrow V=f(X)$ the homomorphism between $X$ and $V$ induced by $f$. Then we have $f=j\circ f'$, where $f'$ is the homeomorphism. 
					
					Then, there exists an isomorphism of sheaves
					\begin{gather*}
						f^\#:f^{-1}\mathcal{O}_Y\rightarrow \mathcal{O}_X
					\end{gather*}
					
					because of open immersion, there exists an isomorphism of rings
					\begin{gather*}
						f^\#:\Gamma(U,f^{-1}\mathcal{O}_Y)\rightarrow \Gamma(U,\mathcal{O}_X)
					\end{gather*}
					
					And we have $f^{-1}\mathcal{O}_Y\cong f'^{-1}\mathcal{O}_V$:
					\begin{gather*}
						\mathcal{O}_V=j^{-1}\mathcal{O}_U\Rightarrow f^{-1}\mathcal{O}_Y=(j\circ f')^{-1}\mathcal{O}_Y=f'^{-1}(j^{-1}\mathcal{O}_Y)=f'^{-1}\mathcal{O}_V
					\end{gather*}
					
					thus we have an isomorphism $\Gamma(U,(f')^{-1}\mathcal{O}_V)\rightarrow \Gamma(U,\mathcal{O}_X)$. Now we've seen an isomorphism
					\begin{gather*}
						f'^{-1}(\mathcal{O}_V)\rightarrow \mathcal{O}_X
					\end{gather*}
					
					Consider $(f',f'^\#:f'^{-1}\mathcal{O}_V\rightarrow \mathcal{O}_x)$. 
					
					Uniqueness: Suppose there exists another $f'':(X,\mathcal{O}_X)\rightarrow (V,\mathcal{O}_V)$, where
					\begin{gather*}
						f'=f''\in Top
					\end{gather*}
					
					 then 
					\begin{center}
						\begin{tikzcd}
							\Gamma(U,f^{-1}\mathcal{O}_Y)\arrow[r,"\cong"]\arrow[d,"\cong"]&\Gamma(U,\mathcal{O}_X)\\
							\Gamma(U,(f')^{-1}\mathcal{O}_V)\arrow[ru,bend left=5,"(f')^\#"]\arrow[ru,bend right=5,"(f'')^\#",swap]
						\end{tikzcd}
					\end{center}
					
					The locally ringed space is unique.
					\end{proof}
					
					\begin{lemma}
						Let $f:X\rightarrow Y$ be a morphism of locally ringed spaces. Let $U\subseteq X$, and $V\subseteq Y$ be open subsets.
						
						Suppose $f(U)\subseteq V$. There exists a unique morphism of locally ringed spaces $f|U:U\rightarrow V$ such that the following diagram is a commutative square of locally ringed spaces
						\begin{center}
							\begin{tikzcd}
								U\arrow[r]\arrow[d,"f|U"]&X\arrow[d,"f"]\\
								V\arrow[r]&Y
							\end{tikzcd}
						\end{center}
					\end{lemma}
					\begin{proof}
						Omitted.
					\end{proof}
					
					There exists a factorization
					\begin{gather*}
						X\rightarrow V\rightarrow Y
					\end{gather*}
					
					where $f|U$ is an open immersion.
			
			\subsubsection{Closed Immersion of Locally Ringed Spaces}	
				\begin{Def}
					Let $i:Z\rightarrow X$ be a morphism of locally ringed spaces. We say that $i$ is a \textbf{closed immersion} if:
					\begin{itemize}
						\item $i$ is a homeomorphism of $Z$ onto a closed subset of $X$;
						\item $\mathcal{O}_X\rightarrow i_*\mathcal{O}_Z$ is surjective; Denote
						\begin{gather*}
							\mathcal{I}:=Ker(\mathcal{O}_X\rightarrow i_*\mathcal{O}_Z)
						\end{gather*}
						\item $\mathcal{I}$ is locally generated by sections.
					\end{itemize}
				\end{Def}
				
				\textbf{Remark} Sometimes, the third condition
				\begin{center}
					$\mathcal{I}$ is locally generated by sections
				\end{center}
				
				is not necessary. (Especially, for the theory of geometry). 
				
				The main problem is 
				\begin{itemize}
					\item Partition of Unity;
					\item Bump functions;
				\end{itemize}
				
				In complex, we have the following theorem
				\begin{theorem}
					(Oka's Coherence Theorem) The sheaf of holomorphic functions $\mathcal{O}^{an}_X$ on complex is a coherent sheaf. 
				\end{theorem}
				
				\begin{example}
					Objects in $C^\infty(\RR)$: Consider
					\begin{gather*}
						\mathcal{I}(U):=\{f\in C^\infty(U)|f(x)=0,\forall x\in U\cap [0,\infty)\}
					\end{gather*}
					
					\begin{itemize}
						\item Homeomorphism: $[0,\infty)$;
						\item Surjectivity: Whitney's extension theorem
						
						\url{https://etd.ohiolink.edu/acprod/odb_etd/ws/send_file/send?accession=ucin1712912147070133&disposition=inline}
						\item Locally generated: Not hold. Consider the neighborhood of $0$.
					\end{itemize}
				\end{example}
				
				\begin{lemma}
					Let $f:Z\rightarrow X$ be a morphism of locally ringed spaces.
					
					In order for $f$ to be a closed immersion, it suffices that there exists an open covering $X=\cup U_i$ such that
					\begin{gather*}
						f^{-1}U_i\rightarrow U_i
					\end{gather*}
					
					is a closed immersion.
				\end{lemma}
				\begin{proof}
					内容...
				\end{proof}
				
				\subsubsection{Sheaf of Ideals}
				Now, we consider te sheaf ideals, as $\mathcal{O}_X$ submodule, which is locally generated by sections.  Now, the support
				\begin{gather*}
					Z:=supp(\mathcal{O}_X/I)=\{x\in X|(\mathcal{O}_X/\mathcal{I})_x\not=0\}
				\end{gather*}
				
				is closed, with inclusion $i:Z\rightarrow X$. There exists a unique sheaf of rings $\mathcal{O}_Z$ on $Z$, with
				\begin{gather*}
					i_*\mathcal{O}_Z:=\mathcal{O}_X/\mathcal{I}
				\end{gather*}
				
				Then, its stalk
				\begin{gather*}
					(i_*\mathcal{O}_Z)=(\mathcal{O}_X/\mathcal{I})=\mathcal{O}_{X,i(z)}/\mathcal{I}_{i(z)}
				\end{gather*}
				
				\subsubsection{Closed Subspace}
				\begin{Def}
					Let $X$ be a locally ringed space, $\mathcal{I}$ be a sheaf of ideals on $X$, which is locally generated by sections.
					
					Then, the \textbf{locally ringed space} is the \textbf{closed subspace} of $X$, associated to the sheaf of ideals $\mathcal{I}$.
				\end{Def}
				
				\begin{lemma}
					Let $f:X\rightarrow Y$ be a closed immersion of locally ringed spaces. Let 
					\begin{gather*}
						\mathcal{I}:=Ker(\mathcal{O}_Y\rightarrow f_*\mathcal{O}_X)
					\end{gather*}
					
					Let $i:Z\rightarrow Y$ be the closed subspace associated to $\mathcal{I}$. There exists a unique isomorphism $f':X\cong Z$ of locally ringed spaces such that
					\begin{gather*}
						f=i\circ f'
					\end{gather*}
				\end{lemma}
				\begin{proof}
					内容...
				\end{proof}
				\begin{lemma}
					Let $X,Y$ be locally ringed spaces. $\mathcal{I}\subseteq \mathcal{O}_X$ be a sheaf of ideals, locally generated by sections.
					
					Let $i:Z\rightarrow X$ be theassociated closed subspace. 
					
					Then, a morphism $f:Y\rightarrow X$ factors through $Z\Leftrightarrow$ the map $f^*\mathcal{I}\rightarrow f^*\mathcal{O}_x=\mathcal{O}_Y$ is zero.
					
					If this is the case the morphism $g:Y\rightarrow Z$ such that $f=i\circ g$ is unique.
				\end{lemma}
				\begin{proof}
					内容...
				\end{proof}
				
				
				\begin{proposition}
					Let $f:X\rightarrow Y$ be a morphism of locally ringed spaces, $\mathcal{I}\subseteq \mathcal{O}_Y$ be a sheaf of ideals, which is locally generated by sections.
					
					Let $i:Z\rightarrow Y$ be the closed subspace associated to the sheaf of ideals $\mathcal{I}$, and
					\begin{gather*}
						\mathcal{J}:=Im(f^*\mathcal{I}\rightarrow f^*\mathcal{O}_Y=\mathcal{O}_X)
					\end{gather*}
					
					Then, $\mathcal{J}$ is ideal locally generated by sections.
					
					Moreover, let $i':Z\rightarrow X$ be the associated closed subspace of $X$, there exists a unique morphism of locally ringed spaces
					\begin{gather*}
						f':Z'\rightarrow Z
					\end{gather*}
					
					such that the following diagram is a commutative locally ringed spaces
					\begin{center}
						\begin{tikzcd}
							Z'\arrow[r,"i'"]\arrow[d,"f'"]&X\arrow[d,"f"]\\
							Z\arrow[r,"i"]&Y
						\end{tikzcd}
					\end{center}
					
					Moreover, this diagram is a \textbf{fibre square} in the category of locally ringed spaces.
				\end{proposition}
				\begin{proof}
					\begin{itemize}
						\item Locally generated:
						\item Factorization:
					\end{itemize}
				\end{proof}
				
				\begin{proposition}
					Given a subsheaf $\mathcal{I}$ of $\mathcal{O}_X$. Let $Z=supp(U\mapsto \mathcal{O}_X(U)/\mathcal{I}(U))$. Then, we have the isomorphism
					\begin{gather*}
						\mathcal{O}_X/\mathcal{I}\cong i_*\mathcal{O}_Z
					\end{gather*}
					if $\mathcal{I}$ is locally generated. 
				\end{proposition}
				\begin{proof}
					Now, for any $x\in X$:
					\begin{itemize}
						\item $x\notin Z\Rightarrow \mathcal{O}_X(U)/\mathcal{I}(U)=0\Rightarrow \mathcal{O}_{X,x}=\mathcal{I}_{X,x}$;
						\item $x\in Z\Rightarrow \mathcal{O}_{Z,x}:=\mathcal{O}_{X,x}/\mathcal{I}_x$
					\end{itemize}
					
					The isomorphism
					\begin{gather*}
						i_*\mathcal{O}_{Z,x}=\mathcal{O}_{X,x}/\mathcal{I}_x
					\end{gather*}
					
					By checking definition.
				\end{proof}
				
				The quotient of stalk is also a local ring. So $(Z,\mathcal{O}_Z)$ is a locally ringed space.
		\section{Affine Schemes}
		\subsection{Modules Sheaf}
			
			\begin{lemma}
				Let $R$ be a ring, $f\in R$:
				
				\begin{itemize}
					\item If $g\in R$ and $D(g)\subseteq D(f)$ then:
					\begin{itemize}
						\item $f$ is invertible in $R_g$;
						\item $g^e=af$ for some $e\geq 1,a\in R$;
						
						\textbf{Remark} Equivalently, $f\in \sqrt{(g)}$;
						
						\item Three is a canonical ring map $R_f\rightarrow R_g$;
						\item There exists a canonical $R_f$-module map $M_f\rightarrow M_g$ for any $R$-module $M$;
					\end{itemize}
					\item (Every Open is Compact) Any open covering $D(f)$ can be refined to a finite open covering of the form
					\begin{gather*}
						D(f)=\cup_{i=1}^n D(g_i)
					\end{gather*}
					\item If $g_1,\dots,g_n\in R$, then $D(f)\subseteq \cup D(g_i)\Leftrightarrow g_i$ generate the unit ideal in $R_f$.
				\end{itemize}
			\end{lemma}
			
			\begin{proof}
				\begin{itemize}
					\item With $D(g)\cong Spec(R_g)$.Thus $f$ is invertible in $R_g$. And $f\cdot \frac{a}{g^d}=1\Leftrightarrow (fg^d-a)\in Ann(g^{e-d})\Rightarrow g^e=af$.
					
					Map $R_f\rightarrow R_g,\frac{b}{f^n}\mapsto \frac{a^nb}{g^{ne}}$, or induced by inclusion $Spec(R_g)\hookrightarrow Spec(R_f)$.
					
					Take the tensor product. 
					\item Omitted. In CA.
					\item Omitted. In CA.
				\end{itemize}
			\end{proof}
			\begin{Def}
				Let $R$ be a ring:
				\begin{itemize}
					\item A \textbf{standard open covering} of $Spec(R)$ is a covering
					\begin{gather*}
						Spec(R)=\cup_{i=1}^n D(f_i);f_i\in R
					\end{gather*}
					\textbf{Remark} Topology of $Spec(R)$ is compact, take the standard open covering.
					\item Suppose $D(f)\subseteq Spec(R)$ is a standard open. A \textbf{standard open covering} of $D(f)$ is a covering
					\begin{gather*}
						D(f)=\cup_{i=1}^n D(g_i);g_i\in R
					\end{gather*}
				\end{itemize}
			\end{Def}
			
			\subsection{Modules Sheaf}
			
			\begin{Def}
				Let $R$ be a ring, and $M$ be an $R$-sheaf. We will define a presheaf $\hat{M}$ on the basis of standard opens.
				
				Suppose that $U\subseteq Spec(R)$ is a standard open. If $f,g\in R$ are such that $D(f)=D(g)$, then there are canonical maps
				\begin{gather*}
					M_f\rightarrow M_g\qquad M_g\rightarrow M_f
				\end{gather*}
				Then, define: For $U=D(f)$,
				\begin{gather*}
					\hat{M}(U)=M_f
				\end{gather*}
				Suppose $D(g)\subseteq D(f)$, then there exists a canonical map
				\begin{gather*}
					\hat{M}(D(f))=M_f\rightarrow M_g=\hat{M}(D(g))
				\end{gather*}
				which is a presheaf of abelian groups.
			\end{Def}
			
			\subsubsection{Stalks}
				Take some $p=x\in Spec(R)$, we have
				\begin{gather*}
					\hat{M}_x=colim_{f\in R,f\notin p} M_f
				\end{gather*}
				With the preorder property: Set $\{f\in R,f\notin p\}$ is preordered by the rule
				\begin{gather*}
					f\geq f'\Leftrightarrow D(f)\subseteq D(f')
				\end{gather*}
				Via definition, we have: For $f_1,f_2\in R-p$, we have
				\begin{gather*}
					f_1f_2\subseteq f_1
				\end{gather*}
				By algebra wew could finally conclude, that
				\begin{gather*}
					\hat{M}_x=colim_{f\in R,f\notin p} M_f=M_p
				\end{gather*}
				
				By the contents of sheaf, there exists a correspondence
				\begin{gather*}
					Mod(\mathcal{O}_X)\Leftrightarrow Mod(\mathcal{O}|B)
				\end{gather*}
				
				where $B$ is a basis. So for the standard opens, $\mathcal{O}_{Spec(R)}$ is unique.
			\subsubsection{Sheaf Property}
			If $D(f)=\cup_{i=1}^n D(g_i)$, then the sheaf condition for this covering is equivalent to the exactness:
			\begin{center}
				\begin{tikzcd}
					0\arrow[r]&M_f\arrow[r]&\oplus M_{g_i}\arrow[r]&M_{g_ig_j}
				\end{tikzcd}
			\end{center}
			Notice that: $D(g_i)=D(fg_i)$. Hence rewritten as
			\begin{center}
				\begin{tikzcd}
					0\arrow[r]&M_f\arrow[r]&\oplus M_{fg_i}\arrow[r]&\oplus M_{fg_ig_j}
				\end{tikzcd}
			\end{center}
			This sequence is exact:\ref{Commutative-Algebra-Rings-Exactness-of-Modules}.
			
			\begin{Def}
				Let $R$ be a ring:
				
				\begin{itemize}
					\item The \textbf{structure sheaf} $\mathcal{O}_{Spec(R)}$ of the spectrum of $R$ is the unique sheaf rings $\mathcal{O}_{Spec(R)}$ which \underline{agrees} with $\hat{R}$ on the basis of standard opens;
					\item The locally ringed space $(Spec(R),\mathcal{O}_{Spec(R)})$ is called the \textbf{spectrum} of $R$ denoted by $Spec(R)$;
					\item The sheaf of $\mathcal{O}_{Spec(R)}$-modules extending $\hat{M}$ to all open subsets on $Spec(R)$ is called the \textbf{sheaf of }$\mathcal{O}_{Spec(R)}$-modules associated to $M$. This sheaf is denoted $\hat{M}$ as well.
				\end{itemize}
			\end{Def}
			
			\begin{lemma}
				Let $R$ be a ring, $M$ be an $R$-module. Let $\hat{M}$ be the sheaf of $\mathcal{O}_{Spec(R)}$-modules associated to $M$:
				
				\begin{itemize}
					\item $\Gamma(Spec(R),\mathcal{O}_{Spec(R)})=R$;  $\Gamma(Spec(R),\tilde{M})=M$;
					\item $\Gamma(D(f),\mathcal{O}_{Spec(R)})=R_f,f\in R$;$\Gamma(D(f),\tilde{M})=M_f,f\in R$;
					\item Restriction maps for $D(g)\subseteq D(f)$
					\begin{gather*}
						R_f\rightarrow R_g\qquad M_f\rightarrow M_g
					\end{gather*}
					\item Stalk: $\mathcal{O}_{Spec(R),x}=R_p,\tilde{M}_{x}=M_p$
					\item 	In particular,
					\begin{gather*}
						R-Mod\rightarrow Mod(\mathcal{O}_{Spec(R)})\qquad M\mapsto \tilde{M}
					\end{gather*}
					is an exact functor.
				\end{itemize}
			\end{lemma}
			\begin{proof}
				\begin{itemize}
					\item Consider $D(1)$;
					\item Definition;
					\item Omitted;
					\item Omitted;
					\item We could prove this stalkwise: Notice that, the localization functor is exact, so the sequence in $Mod(\mathcal{O}_X)$ is exact on each $D(f)$. Thus the functor is exact.
				\end{itemize}
			\end{proof}
			\subsection{The Structure Sheaf}
			\begin{Def}
				Firstly, let $A$ be a ring with $X=Spec(A)$. Define the \textbf{structure sheaf
				\begin{gather*}
					\mathcal{O}_X(U):=\{s:U\rightarrow \sqcup_{p\in U} A_p|s\mbox{ saatisfies (i) (ii)}\}
				\end{gather*}}
				
				with:
				\begin{itemize}
					\item For all $p\in U$, we have $s(p)\in A_p$;
					\item For all $p\in U$, there exists $a,b\in A$ and $V\subseteq U$, with $p\in V\subseteq D(b)$ and
					\begin{gather*}
						s(q)=\frac{a}{b}\in A_q\qquad q\in V
					\end{gather*}
					
					and $\rho_{UV}(s)=s|V$.
					
					\textbf{Remark} $b\in A-q$.
				\end{itemize}
			\end{Def}
			
			\textbf{Remark} There exists a natural ring homomorphism
			\begin{gather*}
				A\rightarrow \mathcal{O}_{Spec(A)}(Spec(A))\\
				a\mapsto (p\mapsto \frac{a}{1}\in A_p)
			\end{gather*}
			
			\begin{Def}
				An \textbf{affine scheme} is a locally ringed space $(X,\mathcal{O}_X)$ if there exists a ring $A$, and an isomorphism 
				
				\begin{gather*}
					(X,\mathcal{O}_X)\stackrel{\sim}{\rightarrow} (Spec(A),\mathcal{O}_{Spec(A)})
				\end{gather*}
			\end{Def}
			\subsection{$AffSch$ and $CRing$ Correspondence}
		
			\begin{lemma}
				\label{algebraic-geometry-morphism-of-stalks-of-locally-ringed-spaces}Let $X$ be a locally ringed space. And $Y=Spec(A)$ be an affine scheme. $f:X\rightarrow Y$ be a morphism of locally ringed spaces.
				
				Given a point $x\in X$ and consider the following ring maps
				\begin{center}
					\begin{tikzcd}
						\Gamma(\mathcal{O}_Y,Y)=A\arrow[r,"f^\#"]&\Gamma(\mathcal{O}_X,X)\arrow[r]&\mathcal{O}_{X,x}
					\end{tikzcd}
				\end{center}
				Then, the preimage of $m_x$ is prime in $A$, denoted by $p$, correspondent to $y\in Spec(A)$. We have
				\begin{gather*}
					f(x)=y
				\end{gather*}
			\end{lemma}
			
			\begin{proof}
				\begin{center}
					\begin{tikzcd}
						\Gamma(X,\mathcal{O}_X)\arrow[r]&\mathcal{O}_{X,x}\\
						\Gamma(Y,\mathcal{O}_Y)\arrow[u]\arrow[r]&\mathcal{O}_{Y,y}=A_p\arrow[u]
					\end{tikzcd}
				\end{center}
				
				Take the inverse image(notice, the right upward arrow is local ring map)
			\end{proof}
			
			\begin{lemma}
				Let $X$ be a locally ringed space, $f\in \Gamma(X,\mathcal{O}_X)$, the set
				\begin{gather*}
					D(f):=\{x\in X|\mbox{image }f\notin m_x\}
				\end{gather*}
				is open. Moreover $f|D(f)$ has an inverse.
			\end{lemma}
			
			\begin{proof}
				Firstly, $x\in X\Leftrightarrow \rho_x(f)\notin m_x\Leftrightarrow \rho_x(f)$ is inversable in $\mathcal{O}_{X,x}\Leftrightarrow $Exists $g\in \Gamma(X,\mathcal{O}_X),\rho_x(f)\rho_x(g)=\rho_x(fg)=id_{\mathcal{O}_{X,x}}\Leftrightarrow \rho_x(fg-id_{\mathcal{O}_X(X)})=0\Leftrightarrow$ Exists $g,U,f|U\circ g|U=id_U$. And notice that $x\in U\subseteq D(f)$, so $D(f)$ is open. 
				
				Moreover, $f|D(f)$ has an inverse: Suppose that $f$ has inverse on $U,V$ with $g,h\in \Gamma(X,\mathcal{O}_X)$:
				
				\begin{gather*}
					g|U\circ f|U=f|U\circ g|U=id_U\qquad h|V\circ f|V=f|V\circ h|V=id_V
				\end{gather*}
				
				Then on $U\cap V$, we have
				
				\begin{gather*}
					g|U\cap V=h|U\cap V
				\end{gather*}
				
				With the pasting property of $\mathcal{O}_X$, there exists some $f^{-1}\in \Gamma(X,\mathcal{O}_X)$ and $f^{-1}|D(f)\circ f|D(f)=f|D(f)\circ f^{-1}|D(f)=id_{D(f)}$.
			\end{proof}
			
			Moreover, in the chapter of invertible modules, we have the following conclusion in general:
			
			\begin{lemma}
				In lemma above, if $X=Spec(A)$ is an affine scheme, then, for $f\in A$,
				
				\begin{gather*}
					D(f)=\{p\in Spec(A)|f\notin p\}
				\end{gather*}
				
				agrees with the definition previously.
			\end{lemma}
			
			\begin{proof}
				$X=Spec(A)x=p\in Spec(A)$, then: $\rho_x:A\rightarrow A_p,t\mapsto\frac{t}{1}$. We have
				
				\begin{gather*}
					\rho_x(t)=\frac{t}{1}\notin m_x=pR_p\Leftrightarrow \frac{t}{1}\in (R-p)R_p
				\end{gather*} 
				
				Suppose $\frac{t}{1}=\frac{a}{b},a,b\in R-p$, then $tb=a$. But $t\in p\Rightarrow a\in p$ contradiction. So $t\notin p$. Conversely also hold, because $(R-p)R_p\cap pR_p=\varnothing$ with inversability. So $\Leftrightarrow f\notin p$. 
				
			\end{proof}
			
			\begin{theorem}
				Let \label{Algebraic-Geometry-Affine-Schemes-Morphisms-from-Locally-Ringed-Space-to-Affine-Scheme}
				\begin{itemize}
					\item $X$ be a locally ringed space;
					\item $Y$ be an affine scheme;
				\end{itemize}
				The map
				\begin{gather*}
					Mor(X,Y)\rightarrow Hom(\Gamma(Y,\mathcal{O}_Y),(X,\mathcal{O}_X))\\
					f\mapsto f^\#
				\end{gather*}
				is bijective.
			\end{theorem}
			
			\begin{proof}
				$Y$ is affine. Let $Y=Spec(R)$, then 
				
				$\rightarrow$ Omitted
				
				$\leftarrow$ Construction of the inverse map: Given some 
				\begin{gather*}
					\psi_Y:\Gamma(Y,\mathcal{O}_Y)=R\rightarrow \Gamma(X,\mathcal{O}_X)
				\end{gather*}
				Then, we will construct the corresponding map $\Psi:X\rightarrow Y$. Given $x\in X$, then define
				\begin{center}
					$\Psi(x):=Ker($\begin{tikzcd}
						\Gamma(Y,\mathcal{O}_Y)\arrow[r,"\psi_Y"]&\Gamma(X,\mathcal{O}_X)\arrow[r]&\mathcal{O}_{X,x}\arrow[r]&\mathcal{O}_{X,x}/m_x
					\end{tikzcd}$)$
				\end{center}
				\begin{itemize}
					\item Continuity: Take standard open $D(g)\subseteq Y=Spec(R)$. Then, 
					\begin{gather*}
						\Psi^{-1}(D(g))=D(\Psi_Y(g))
					\end{gather*}
					which is open.
					
					\item Sheaf Structure Homomorphism: We should construct the $\Psi$-map 
					\begin{itemize}
						\item Given by:
						\begin{gather*}
							\Gamma(D(g),\mathcal{O}_Y)=\Gamma(Y,\mathcal{O}_Y)_g\rightarrow \Gamma(\Psi^{-1}(D(g)),\mathcal{O}_X))
						\end{gather*}
						
						And $\Psi_Y(g)$ is invertible on $D(\Psi_Y(g))$. Thus we could extend $\Psi_Y$ to 
						\begin{gather*}
							\psi_{D(g)}:\Gamma(Y,\mathbb{O}_Y)_g\rightarrow \Gamma(\Psi^{-1}(D(g),\mathcal{O}_X))
						\end{gather*} 
						\item The morphism is local for stalks. Suppose the morphism is defined as
						\begin{gather*}
							(\Psi,\psi):(X,\mathcal{O}_X)\rightarrow (Y,\mathcal{O}_Y)
						\end{gather*} 
						
						Now, 
						\begin{gather*}
							\psi_x:\mathcal{O}_{Y,\Psi(x)}\rightarrow \mathcal{O}_{X,x}
						\end{gather*}
						
						it is local: \ref{algebraic-geometry-morphism-of-stalks-of-locally-ringed-spaces}
					\end{itemize}
					\item Finally, theere exists an inverse between $(\Psi,\psi)\mapsto \psi_Y$ and the construction $\psi_Y\mapsto (\Psi,\psi)$. 
				\end{itemize}
			\end{proof}
			
			\begin{corollary}
				The category of affine schemes is equivalent to the opposite of the category of rings.
				
				The equivalence is given by the functor that associates to an affine scheme the global sections of its structure sheaf.
			\end{corollary}
			\begin{proof}
				Omitted.
			\end{proof}
			
			\begin{corollary}
				$Spec(\ZZ)$ is the final object in the category of affine scheme/locally ringed space.\label{algebraic-geometry-affine-scheme-specZ-is-final-object-of-locally-ringed-space}
			\end{corollary}
			\begin{proof}
				Suppose that $(X,\mathcal{O}_X)$ be a locally ringed space. Then, we have $Hom(X,Spec(\ZZ))\cong Hom(\Gamma(Spec(\ZZ),Spec(\ZZ)),\Gamma(X,\mathcal{O}_X))=Hom(\ZZ,\mathcal{O}_X(X))$. And the morphism $\ZZ\rightarrow \mathcal{O}_X(X)$ is well-defined.
			\end{proof}
			\subsection{Product, Fibre Product}
			\begin{lemma}
				The category of affine schemes has:
				\begin{itemize}
					\item Finite products;
					\item Fibre products;
				\end{itemize}
				In other words, it has finite limits. 
				
				Moreover, the products and fibre products in the category of affine schemes are the same as in the category of locally ringed spaces. In formula
				\begin{gather*}
					Spec(R)\times Spec(S)=Spec(R\otimes_\ZZ S)\\
					Spec(A)\times_{Spec R} Spec(B)=Spec(A\otimes_R B)
				\end{gather*}
			\end{lemma}
			\begin{proof}
				\begin{itemize}
					\item 
					\item 
					\item 
				\end{itemize}
			\end{proof}
			
			\subsection{Disjointed Union}
			
			\begin{lemma}
				Let $X$ be a locally ringed space. Assume that:
				\begin{gather*}
					X=U\sqcup V
				\end{gather*}
				where $U,V$ are affine opens. Then $X$ is an affine scheme.
			\end{lemma}
			
			\begin{proof}
				Define: 
				\begin{gather*}
					R=\mathcal{O}_X(U)\times \mathcal{O}_X(V)
				\end{gather*}
				Then, there exists a canonical morphism of locally ringed spaces
				\begin{gather*}
					内容...
				\end{gather*}
			\end{proof}
		
		\subsection{Quasi-Coherent Sheaf on Affines}
			We will prove the equivalence of following categories
			\begin{gather*}
				QCoh(\mathcal{O}_{Spec(R)})\Leftrightarrow Mod-R
			\end{gather*}
			
			Where $QCoh(\mathbb{O}_X)$ is a subcategory of $Mod(\mathcal{O}_X)$.
			
			\begin{lemma}
				Let: 
				\begin{itemize}
					\item $(X,\mathcal{O}_X)=(Spec(R),\mathcal{O}_{Spec(R)})$ be an affine scheme; 
					\item $M$ be an $R$-module. 
				\end{itemize}
				
				Then, there exists an isomorphsim between sheaf $\hat{M}$ associated to the R-module $M$ and the sheaf $\mathcal{F}_M$ associated to the R-module $M$. This isomorphism is functorial in module $M$. 
				
				Moreover they are characterized by the following mapping property
				\begin{gather*}
					Hom_{\mathcal{O}_X}(\hat{M},\mathcal{F})=Hom_R(M,\Gamma(X,\mathcal{F}))
				\end{gather*}
				for any sheaf of $\mathcal{O}_X$-modules $\mathcal{F}$. Here a map $\alpha:\hat{M}\rightarrow \mathcal{F}$ corresponds to its effect on global sections.
			\end{lemma}
			\begin{proof}
				内容...
			\end{proof}
			
			\begin{lemma}
				Let $(X,\mathcal{O}_X)=(Spec(R),\mathcal{O}_{Spec(R)})$ be an affine scheme. There exists canonical isomorphisms
				\begin{itemize}
					\item $\widetilde{M\otimes_R N}=\widetilde{M}\otimes_{\mathcal{O}_{Spec(R)}}\widetilde{N}$;
					\item $\widetilde{T^nM}\cong T^n(\widetilde{M}),\widetilde{Sym^n(M)}\cong Sym^n(\widetilde{M}),\widetilde{\wedge^n(M)}=\wedge^n (\widetilde{M})$;
					\item If $M$ is finitely presented, then
					\begin{gather*}
						Hom_{\mathcal{O}_X}(\tilde{M},\tilde{N})\cong \widetilde{Hom_R(M,N)}
					\end{gather*}
				\end{itemize}
			\end{lemma}
			\begin{proof}
				内容...
			\end{proof}
			
			\begin{lemma}
				内容...
			\end{lemma}
			\begin{proof}
				内容...
			\end{proof}
			
			\begin{lemma}
				内容...
			\end{lemma}
			\begin{proof}
				内容...
			\end{proof}
			
			\begin{theorem}
				Let $(X,\mathcal{O}_X)=(Spec(R),\mathcal{O}_{Spec(R)})$ be affine scheme. The functors
				\begin{gather*}
					M\mapsto \tilde{M}\qquad \mathcal{F}\mapsto \Gamma(X,\mathcal{F})
				\end{gather*}
				give a rise to the equivalence of functors
				\begin{center}
					\begin{tikzcd}
						Mod(\mathcal{O}_X)\arrow[r,shift left=1]&Mod-R\arrow[l,shift left=1]
					\end{tikzcd}
				\end{center}
			\end{theorem}
			\begin{proof}
				内容...
			\end{proof}
			\subsection{Kernels,Cokernels}
			\begin{lemma}
				Let $X=Spec(R)$ be an affine scheme. Kernels and cokernels of maps of quasi-coherent $\mathcal{O}_X$-modules are quasi-coherent.
			\end{lemma}
			\begin{proof}
				
			\end{proof}
			\subsection{Direct Sum}
			\begin{lemma}
				内容...
			\end{lemma}
			\begin{proof}
				内容...
			\end{proof}
			\subsection{Exact Sequence}
				\begin{lemma}
				Let $(X,\mathcal{O}_X)=(Spec(R),\mathcal{O}_{Spec(R)})$ be an affine scheme. Suppose that
				\begin{gather*}
					0\rightarrow\mathcal{F}_1\rightarrow \mathcal{F}_2\rightarrow \mathcal{F}_3\rightarrow 0
				\end{gather*}
				is a short exact sequence of sheaves of $\mathcal{O}_X$-modules. If two out of three are quasi-coherent so is the third.
			\end{lemma}
			\begin{proof}
				内容...
			\end{proof}
			\subsection{Open and Closed Subspaces of Affines}
			
			In general, the open/closed subspaces and immersions are discussed in locally ringed spaces.
			
			\subsubsection{Open Subspace}
				\begin{lemma}
				Let $Y$ be an affine scheme, $f\in \Gamma(Y,\mathcal{O}_Y)$, the open subspace $D(f)$ is an affine subspace. 
			\end{lemma}
			\begin{proof}
				As a matter of fact, $D(f)$ has a structure of ringed space by
				\begin{gather*}
					\Gamma(D(f),\mathcal{O}_{D(f)})=\Gamma(Y,\mathcal{O}_Y)_f\\
					\phi:U=Spec(R_f)\rightarrow Spec(R)=Y
				\end{gather*}
				
				We have know:
				\begin{itemize}
					\item $Spec(R_f)\rightarrow Spec(R)$ is a homeomorphism onto $D(f)$;
					\item $\phi^{-1}\mathcal{O}_Y\rightarrow \mathcal{O}_U$ is an isomorphism on stalk. Thus is an open immersion because isomorphism of sheaves.
				\end{itemize}
			\end{proof}
			\subsubsection{Closed Subspace}
			\begin{example}
				Let: $R$ be a ring, $I\subseteq R$ be an ideal. Consider the morphism of affine schemes
				\begin{gather*}
					i:Z=Spec(R/I)\rightarrow Spec(R)=X
				\end{gather*}
				
				This is a homeomorphism of $Z$ onto the closed subset of $X$. 
				
				Moreover, if $I\subseteq p\subseteq R$ is a prime corresponding ot a point $x=i(z)$, where $x\in X,z\in Z$, then on stalks we get the map
				\begin{gather*}
					\mathcal{O}_{X,x}=R_p\rightarrow R_p/IR_p=\mathcal{O}_{Z,z}
				\end{gather*} 
				Thus we see that $i$ is a closed immersion of locally ringed spaces.
			\end{example}
			\begin{lemma}
				Let $(X,\mathcal{O}_X)=(Spec(R),\mathcal{O}_{Spec(R)})$ be an affine scheme. Let $i:Z\rightarrow X$ be any closed immersion of locally ringed spaces. Then, there exists a unique ideal $I\subseteq R$ such that:$i:Z\rightarrow X$ can be identified with the closed immersion
				\begin{gather*}
					Spec(R/I)\rightarrow Spec(R)
				\end{gather*}
				constructed above.
			\end{lemma}
			\begin{proof}
				Such $I$ is given by: The ideal sheaf $\mathcal{I}=Ker(\mathcal{O}_X\rightarrow i_*\mathcal{O}_Z)$ on $Spec(R)$. $I=\Gamma(Spec(R),\mathcal{I})$ is an ideal.(By kernel)
			\end{proof}
			
			And such $\hat{I}$ is naturally locally generated...
		\section{Topological Properties}
			More discussion should be listed in CA.
	\chapter{Schemes}
		\section{Definitions}
			
		\begin{Def}
			A \textbf{scheme} is a locally ringed space $(X,\mathcal{O}_X)$, with
			
			\begin{itemize}
				\item Open Covering $X=\cup U_i$;
				\item $(U_i,\mathcal{O}_X|U_i)$ is affine scheme.
			\end{itemize}
		\end{Def}
		
		\subsection{Morphisms of Schemes}
		
		\begin{Def}
			The morphism of schemes is the morphism of locally ringed spaces.
		\end{Def}
		\subsubsection{Affinization Morphism}
		\begin{Def}
			\label{Algebraic-Geometry-Affinization-Map}Define the functor:
			\begin{gather*}
				aff:Sch\rightarrow Affsch\qquad X\mapsto Spec(\mathcal{O}_X(X))
			\end{gather*}
			And induces a natural morphism $X\rightarrow Spec(\mathcal{O}_X(X))$
		\end{Def}
		Especially,
		\begin{gather*}
			(aff_X,aff_X^\#):X\rightarrow Spec(\mathcal{O}_X(X))
		\end{gather*}
		\begin{itemize}
			\item $\phi(x)=aff_X(x):=[\mathcal{O}_X(X)\rightarrow \mathcal{O}_{X,x}]^{-1}(m_x)$
			\item $aff_X(D(f)):\mathcal{O}_X(X)_f\rightarrow \Gamma(\phi^{-1}(D(f)),\mathcal{O}_X)$ induced by the morphism $\mathcal{O}_X(X)\rightarrow$. 
			
			$f|D(f))$ is inversable. 
		\end{itemize}
		Specified by the module homomorphism
		\begin{gather*}
			Hom(X,Spec(\mathcal{O}_X(X)))=Hom(\Gamma(Spec\mathcal{O}_X(X),\mathcal{O}_{Spec(\mathcal{O}_X(X))}),\Gamma(X,\mathcal{O}_X))\\
			\qquad=Hom(\mathcal{O}_X(X),\mathcal{O}_X(X))
		\end{gather*}
		Take the identity.
		
		
		\begin{lemma}
			Firstly, $aff$ is a covariant functor.
		\end{lemma}
		\begin{proof}
			Consider morphism $\varphi\in Hom(X,Y)$. Then, we have $aff(\varphi):Spec(\mathcal{O}_Y(Y))\rightarrow Spec(\mathcal{O}_X(X))$. 
		\end{proof}
		\begin{proposition}
			The morphism $X\rightarrow Spec(A)$ factors through $aff_X$
			\begin{gather*}
				X\rightarrow Spec(\mathcal{O}_X(X))\rightarrow Spec(A)
			\end{gather*}
		\end{proposition}
		\begin{proof}
			Notice that: There exists $f^\#(Spec(A)):A\rightarrow \mathcal{O}_X(X)$. We have correspondence $aff_X\Leftrightarrow id_{\mathcal{O}_X(X)}$. Then the factorization $Spec(\mathcal{O}_X(X))\rightarrow Spec(A)$ is $Spec(f^\#(Spec(A)))$. Then consider the composition of $f^\#(Spec(A))\circ aff_X^\#=f^\#$
		\end{proof}
		\begin{proposition}
			Scheme $X$ is affine iff $aff_X$ is an isomorphism.
		\end{proposition}
		\begin{proof}
			Omitted.
		\end{proof}
		\subsubsection{Open Immersion}
		
		\begin{lemma}
			Let $X$ be a scheme, $j:U\rightarrow X$ be an open immersion of locally ringed spaces.
			
			Then, $U$ is a scheme. In particular, any open subspace of $X$ is a scheme.
		\end{lemma}
		
		\begin{proof}
			Suppose $u\in U$. Then, select affine neighborhood $u\in V$, there must be some $u\in D(f)\subseteq U$.
			
			$D(f)$ is affine scheme, so $U$ is a scheme via $U=\cup_{u\in U} D(f_u)$, with paasting of sheaves $\mathcal{O}_X|D(f_u)$.
		\end{proof}
		
		\begin{corollary}
			Any scheme $X$ has a basis for topology consisting of affine opens $D(f)$ where $f\in \Gamma(U,\mathcal{O}_X|U)=(Spec(R),\mathcal{O}_{Spec(R)})$.
		\end{corollary}
		\begin{proof}
			Omitted.
		\end{proof}
		\subsubsection{Closed Immersion}
		\begin{lemma}
			Let $X$ be a scheme, $i:Z\rightarrow X$ be a closed immersion of locally ringed spaces:
			\begin{itemize}
				\item The locally ringed space $Z$ is a scheme;
				\item The kernel $\mathcal{I}$ of the map $\mathcal{O}_X\rightarrow i_*\mathcal{O}_Z$ is a \textbf{quasi-coherent sheaf};
				\item For any affine open $U=Spec(R)$ of $X$, the morphism
				\begin{gather*}
					i^{-1}(U)\rightarrow U\qquad Spec(R/I)\rightarrow Spec(R)
				\end{gather*}
				for some ideal $I\subseteq R$ and
				We have $\mathcal{I}|U=\hat{I}$.
			\end{itemize}
		\end{lemma}
		\begin{proof}
			
		\end{proof}
		
		\begin{lemma}
			Let $f:Y\rightarrow X$ be an immersion of schemes. Then $f$ is closed immersion iff $f(Y)\subseteq X$ is a closed subset.(immersion: Here \ref{Algebraic-Geometry-Schemes-Immersion})
		\end{lemma}
		\begin{proof}
			$\Rightarrow$ Omitted.
			
			$\Leftarrow$ 
		\end{proof}
		\subsubsection{Immersion}
		\begin{Def}
			\begin{itemize}
				\item The morphism of schemes is called an \textbf{open immersion} if it is an open immersion of locally ringed spaces;
				\item An \textbf{open subscheme} of $X$ is an open subspace of $X$ is an open subspace of $X$ or a scheme by taking subset;
				\item A morphism of schemess is called a \textbf{closed immersion} if it is a closed immersion of locally ringed spaces;
				\item A \textbf{closed subscheme} of $X$ is a closed subspace of locally ringed space $X$ or by lemma above;
				\item An \textbf{immersion} can be factored as $f=j\circ i$ where $i$ is closed immerions, $j$ is open immersion.\label{Algebraic-Geometry-Schemes-Immersion}
			\end{itemize}
		\end{Def}
		
			More contents of morphisms of schemes will be given in another chapter.
		\subsection{Zariski Topology of Schemes}
		
			\begin{lemma}
				Let $X$ be a scheme. A \textbf{irreducible closed subset} of $X$ has a unique generic point.
				
				In other words, it is a sober space.
			\end{lemma}
			
			\begin{proof}
				Select irrducible set $Z$
				
				Existence of generic point: For every affine open $U=Spec(R)$ and we know that
				\begin{gather*}
					Z\cap U=V(I)
				\end{gather*}
				
				for a unique radical ideal $I\subseteq R$. And $V(I)$ is irreducible$\Leftrightarrow$ $I$ is prime. Thus $p=I$ is the generic point of $Z\cap U$ by $Z\cap U=V(I)$.
				
				Now, suppose $\xi'$ be another generic point, then there exists some neighborhood $U$ contains $\xi,\xi'$ by definition. Finally $\xi'=\xi$. 
			\end{proof}
			
			\begin{lemma}
				The underlying topological space of any shceme is locally quasi-compact.
			\end{lemma}
			\begin{proof}
				Affine open is compact. Omitted.
			\end{proof}
			\begin{lemma}
				Let $X$ be a scheme, $U,V$ be affine opens of $X$, $x\in U\cap V$.
				
				There exists an affine open neighborhood $W$ of $x$ such that $W$ is a standard open of both $U$ and $V$.
			\end{lemma}
			\begin{proof}
				Let $U=Spec(R)$ and $V=Spec(S)$. Then, suppose $p=q=x\in U\cap V;p\subseteq R,q\subseteq S$. 
			\end{proof}
			\begin{lemma}
				Let $X$ be a scheme and $X=\cup U_i$ be an affine open covering.Then, if $V$ is an affine open, there exists a standard open covering $V=\cup V_i$ where $V_j\subseteq U_i$ is a standard open.
			\end{lemma}
			\begin{proof}
				内容...
			\end{proof}
			\begin{lemma}
				Let $X$ be a scheme and $\mathcal{B}$ be the set of affine opens of $X$. Let $\mathcal{F}$ be a presheaf of sets on $\mathcal{B}$. The followings are equivalent:
				\begin{itemize}
					\item $\mathcal{F}$ is the restriction of sheaf on $X$ to $\mathcal{B}$;
					\item $\mathcal{F}$ on $\mathcal{B}$ and
					\item $\mathcal{F}(\varnothing)$ is a singeleton;
					\item Whenever $U=V\cup W$ with $U,V,W\in \mathcal{B}$ and $V,W\subseteq U$ standard open. the map 
					\begin{gather*}
						\mathcal{F}(U)\rightarrow \mathcal{F}(V)\times \mathcal{F}(W)
					\end{gather*} 
					is injetive with image the set of pair $(s,t)$ such that $s|V\cap W=t|V\cap W$ 
				\end{itemize}
			\end{lemma}
			\begin{proof}
				内容...
			\end{proof}
			
			\subsubsection{Special Case:Scheme on Finite discrete Set}
			\begin{lemma}
				Let $X$ be a finite and discrete topological space and $(X,\mathcal{O}_X)$ be a scheme. Then it is affine.
			\end{lemma}
			\begin{proof}
				Now consider $\coprod \{x_i\}$, then each $x_i$ must be affine. Thus $X$ is the finite disjoint union of affine schemes. Hence is affine by taking disjoint union.
			\end{proof}
			\begin{example}
				Existence of a scheme without closed point:
			\end{example}
			
			\subsection{Points in Schemes}\label{Algebraic-Geometry-Points-in-Schemes}
			\begin{Def}
				Given a scheme $X$ we can define a functor
				\begin{gather*}
					h_X:Sch^{opp}\rightarrow Sets\qquad T\mapsto Mor(T,X)
				\end{gather*}
				This is called the \textbf{functor of points} of $X$. 
			\end{Def}
			Let 
			\begin{itemize}
				\item $X$ be a scheme;
				\item $R$ be a local ring with maximal ideal $m$;
			\end{itemize}
			Suppose that $f:Spec(R)\rightarrow X$ be a morphisms of schemes. Then, $x:=f(m)$, we obtain a local homomorphsim of local rings
			\begin{gather*}
				f^\#:\mathcal{O}_{X,x}\rightarrow \mathcal{O}_{Spec(R),m}=R_m=R
			\end{gather*}
			
			\textbf{Remark} $R_m=R$. For local ring $R$, $R-m$ is the collection of inversable elements. Thus,
			\begin{lemma}
				Let $X$ be a scheme, $R$ be a local ring. Then, the construction above gives a bijective correspondence
				\begin{center}
					Morphisms $Spec(R)\rightarrow X\Leftrightarrow$ Pairs $(x,\varphi)$ consisting of a point $X$ and a local homomorphism of local rings $\mathcal{O}_{X,x}\rightarrow R$. 
				\end{center}
			\end{lemma}
			\begin{proof}
				Firstly, we should consider the affine case: Let $X=Spec(A)$. Then, $\psi:A\rightarrow R$ has a factorization of local rings
				\begin{gather*}
					A\rightarrow A_p\rightarrow R
				\end{gather*}
				where $p=\psi^{-1}(m)$. Thus $Spec(R)\rightarrow X$ is given by composition of $A\rightarrow A_p$ and $A_p\rightarrow R$, which is the local homomorphism $\mathcal{O}_{X,x}\rightarrow R$.
				
				Now consider the general case. $\Leftarrow$ Consider $\mathcal{O}_X(U)\rightarrow \mathcal{O}_{X,x}\rightarrow R,Spec(R)\rightarrow U$ where $U$ is affine. And paste the morphisms for sheaf. 
				
				Conversely, given $f:Spec(R)\rightarrow X$, we must prove: For \underline{any} affine open of $X$ such that $x=f(m)\in U$, $Im(f)\subseteq U$. In fact: $m\in f^{-1}(U)$, so $f^{-1}(U)=Spec(R)$.
			\end{proof}
			As a special case of the lemma above we obtain for every point $x$ of a scheme $X$ a canonical morphism
			\begin{gather*}
				内容...
			\end{gather*}
			
			\begin{center}
				\begin{tikzcd}
					Spec(\mathcal{O}_{X,x})\arrow[r]\arrow[d]&X\arrow[d]\\
					Spec(\mathcal{O}_{S,s})\arrow[r]&S
				\end{tikzcd}
			\end{center}
			\subsubsection{Generalization,Specialization}
			
			\begin{Def}
				\begin{itemize}
					\item Let $x,x'\in X$. If $x'\in \overline{\{x\}}$, we say that 
					
					\begin{itemize}
						\item $x'$ is a specialization of $x$;
						\item $x$ is a generalization of $x'$
					\end{itemize}
					\item A subset $V\subseteq X$ is called \textbf{closed under specialization/generalization} if it contains all specializations(generalizations) of all its points.
 				\end{itemize}
			\end{Def}
			
			\begin{lemma}
				Let $X$ be a scheme, $x,x'\in X$. Then:
				\begin{center}
					$x'\in X$ is a generalization of $x\Leftrightarrow x'$ is in the image of canonical morphism $Spec(\mathcal{O}_{X,x})\rightarrow X$. 
				\end{center}
			\end{lemma}
			\begin{proof}
				内容...
			\end{proof}
			
			\begin{example}
				内容...
			\end{example}
			\subsubsection{Points-Spectra Field Correspondence}
			\begin{lemma}
				Let $X$ be a 
			\end{lemma}
			\begin{proof}
				内容...
			\end{proof}
			
			\begin{Def}
				Let $f:X\rightarrow Y$ be a continuous map, we say that
				
				\begin{itemize}
					\item 
					\item 
				\end{itemize}
			\end{Def}
			\subsubsection{Valuative Criterion for Universal Closedness}
			Review: The \textbf{closed map} is the map which maps closed sets to closed sets.
			\begin{Def}
				
			\end{Def}
			
			\subsubsection{Valuative Criterion of Separatedness}
		
		\subsection{Representability}
			\begin{Def}
				Suppose that $F:Sch^{op}\rightarrow Sets$ be functor. Then, $F$ is \textbf{representable} if there exists a scheme $X$ such that $h_X\cong F$.
			\end{Def}
			Now apply \underline{Yoneda's lemma} we have
			\begin{gather*}
				Hom(h_X,F)\cong F(X)
			\end{gather*}
			Thus $(X,s:h_X\rightarrow F)$ is unique up to unique isomorphism if it exists.
			
			Moreover, we have a bijection
			\begin{gather*}
				Mor_{Fun(Sch^{op},Sets)}(h_Y,F)\rightarrow F(Y)\qquad s\mapsto s(id_Y)
			\end{gather*}
			Conversely, we have
			\begin{gather*}
				F(Y)\rightarrow Mor_{Fun(Sch^{op},Sets)}(h_Y,F)\qquad \xi\mapsto s_\xi
			\end{gather*}
			where $s_\xi:h_Y\rightarrow F,h_Y(T)=Mor(T,Y)\rightarrow F(T),f\mapsto f^*\xi$
			
			\textbf{Remark} In particular, in hte case that $F$ is representable, there exists a scheme $X$ and an element $\xi\in F(X)$ We say
			\begin{Def}
				Pair $(X,\xi)$ represents $F$, if $F=s_\xi$.
			\end{Def}
			
			\begin{Def}
				Such $\xi\in F(X)$ is often called the \textbf{universal family} for reasons that will become more clear when we talk about [algebraic stacks].
			\end{Def}
			
			Then, every element $\xi'\in F(T)$ for any $T$ is of the form $\xi'\in F(T)$ for any $T$ is of the form $\xi'=f^*\xi$ for a unique morphism $f:T\rightarrow T$.
			
			\begin{example}
				Consider the rule which associates to every scheme $T$ the set
				\begin{gather*}
					F(T)=\Gamma(T,\mathcal{O}_T)
				\end{gather*}
				Then, we can trun this into a contravariant functor by using for a morphism $f:T\rightarrow T'$ the pullback $f^\#:\Gamma(T,\mathcal{O}_{T})\rightarrow \Gamma(T',\mathcal{O}_{T'})$. 
				
				Thus, given a ring $R$ with element $t\in R$ there exists a unique ring homomorphism $\ZZ[x]\rightarrow R$ which maps $x$ to $t$, thus we have a bijection $R\cong Hom(\ZZ[x],R)$,
				\begin{gather*}
					Mor(T,Spec(\ZZ[x]))=Hom(\ZZ[x],\Gamma(T,\mathcal{O}_T))=\Gamma(T,\mathcal{O}_T)
				\end{gather*}
				This give an isomorphism of functors
				\begin{gather*}
					F\rightarrow h_{Spec(\ZZ[x])}
				\end{gather*}
			\end{example}
			
			
			
			\subsubsection{Criterion for Representability}
			
			\begin{Def} Let $F$ be a contravariant functor on category of schemes with values in sets:
				\begin{itemize}
					\item We say $F$ \underline{satisfies the sheaf property for the Zariski topology} if for every scheme $T$ and every open covering $T=\cup_{i\in I} U_i$ and for any collection of elements $\xi_i\in F(U_i)$ such that
					\begin{gather*}
						\xi_i|U_i\cap U_j=\xi_j|U_i\cap U_j
					\end{gather*}
					There exists a unique element $\xi\in F(T)$ such that $\xi_i=\xi|U_i\in \mathcal{F}(U_i)$;
					\item A \textbf{subfunctor } $H\subseteq T$ is a rule associate to every scheme $T$ a subset
					\begin{gather*}
						H(T)\subseteq F(T)
					\end{gather*} 
					And for $T\rightarrow T'$ we have $F(f):F(T)\rightarrow F(T');H(f)=F(f)|H(T):H(T)\rightarrow H(T')$.
					\item Let $H\subseteq F$ be a subfunctor. We say $H\subseteq F$ is \textbf{representable by open immersions} if for all pairs $(T,\xi)$, where $T$ is a scheme and $\xi\in F(T)$ there exists an open subscheme $U_\xi\subseteq T$ with following properties
					\begin{center}
						A morphism $f:T'\rightarrow T$ factors through $U_\xi$ iff $f^*\xi\in H(T')$.
					\end{center}
					\item Let $I$ be a set. For each $i\in I$ $H_i\subseteq F$ be a subfuncotr. We say that the collection $(H_i)_{i\in I}$ covers $F$ iff for every $\xi\in F(T)$ there exists an open covering $T=\cup U_i$ such that $\xi|U_i\in H_i(U_i)$
				\end{itemize}
			\end{Def}
			
			\begin{lemma}
				Let $F$ bea covariant functor on the category of schemes with value in the category of sets. Suppose that:
				\begin{itemize}
					\item $F$ satisfies the sheaf property for the Zariski topology;
					\item There exists a set $I$ and a collection of subfunctors such that:
					\begin{itemize}
						\item each $F_i$ is representable;
						\item each$F_i\subseteq F$ is representable by open immersions;
						\item the collection $(F_i)_{i\in I}$ covers $F$
					\end{itemize}
					Then, $F$ is representable.
				\end{itemize}
			\end{lemma}
			\begin{proof}
				内容...
			\end{proof}
			
			\textbf{Remark} Suppose the functor $F$ is defined on all locally ringed spaces, and if conditions of above are replaced by the following:
			\begin{itemize}
				\item $F$ satisfies the sheaf property on the category of locally ringed spaces;
				\item There exists a set $I$ and a collection of subfunctors $F_i\subseteq F$ such that:
				\begin{itemize}
					\item Each $F_i$ is representable by a scheme;
					\item Each $F_i\subseteq F$ is representable by open immersions on the category of locally ringed spaces; and
					\item The collection $(F_i)_{i\in I}$ cover $F$ as a functor on the category of locally ringed spaces.
				\end{itemize}
			\end{itemize}
		\section{Classification of Schemes(by Rings)}
			
			\subsection{Reduced Scheme}
			
				
				\subsection{Integral, Irreducible and Reduced Scheme}
					
					\subsubsection{Integral Scheme}
						\begin{Def}
							Let $X$ be a scheme. Then $X$ is \textbf{integral} if for every nonempty affine open $Spec(R)=U$, $R$ is an integral domain.
						\end{Def}
						
					\subsubsection{Irreducible Scheme}
						\begin{Def}
							Scheme $X$ is \textbf{irreducible} iff its underlying topological space $|X|$ is irreducible.
						\end{Def}
						
						\begin{lemma}
							The followings are equivalent:
							\begin{itemize}
								\item Scheme $X$ is irreducible;
								\item There exists an affine open covering $X=\cup_i U_i$ such that $I$ is not empty, $U_i$ is irreducible for all $i\in I$. And $U_i\cap U_j\not=\varnothing$ for all $i,j\in I$; 
								\item The scheme $X$ is nonempty and every nonempty affine open $U\subseteq X$ is irreducible;
							\end{itemize}
						\end{lemma}
						\begin{proof}
							内容...
						\end{proof}
					\subsubsection{Integral$\Leftrightarrow$ Reduced+Irreducible}
						\begin{proposition}
							A scheme $X$ is integral iff it is reduced and irreducible.
						\end{proposition}
						\begin{proof}
							$\Rightarrow$
							
							$\Leftarrow$
						\end{proof}
				\subsection{Types of Schemes defined by Propreties of Rings}
					\subsubsection{Local Properties}
						\begin{Def}
							\begin{itemize}
								\item Let $P$ be a \underline{local property of rings}. We say $P$ is local if the following hold:
								\begin{itemize}
									\item 
									\item 
								\end{itemize}
								\item Let $X$ be a scheme. We say $X$ is locally $P$ if for any $x\in X$ there exists an affine open neighborhood $U$ of $x$ in $X$ such that $\mathcal{O}_X(U)$ has property $P$.
							\end{itemize}
						\end{Def}
						\begin{lemma}
							内容...
						\end{lemma}
						\begin{proof}
							内容...
						\end{proof}
						
						\begin{lemma}
							内容...
						\end{lemma}
						\begin{proof}
							内容...
						\end{proof}
						
						\begin{lemma}
							The following properties of a ring $R$ are local:
							\begin{itemize}
								\item 
								\item 
								\item 
								\item 
							\end{itemize}
						\end{lemma}
						\begin{proof}
							内容...
						\end{proof}
				\subsection{Noetherian Schemes}
				
				\subsection{Jacobson Schemes}
					\begin{Def}
						内容...
					\end{Def}
				
				\subsection{Normal Schemes}
					\begin{Def}
						A scheme $X$ is normal iff for all $x\in X$ the local ring $\mathcal{O}_{X,x}$ is a normal domain.
					\end{Def}
				
			
					
			\section{Fibre Products of Schemes}
				\subsection{Schemes as Functors}
					Here: \ref{Algebraic-Geometry-Points-in-Schemes}
				\subsection{Fibre Products}
			
				\subsubsection{Existence of Fibre Products of Schemes}
				
					\begin{lemma}
						The category of schemes has a
						
						\begin{itemize}
							\item Final Objects;
							\item Products;
							\item Fibre Products;
						\end{itemize}
						
						In other words, the category of schemes has finite limits.
					\end{lemma}
					
					\begin{proof}
						Review: \ref{algebraic-geometry-affine-scheme-specZ-is-final-object-of-locally-ringed-space}
						
						Let $f:X\rightarrow S,g:Y\rightarrow S$ be morphisms of schemes. We have to show that the functor
						\begin{gather*}
							F:Sch\rightarrow Sets\qquad T\mapsto Mor(T,X)\times_{Mor(T,S)} Mor(T,Y)
						\end{gather*}
						
						is representable. We should apply the lemma: 
					\end{proof}
					Another proof through affine case:
					\begin{lemma}
						(Affine Case)Suppose that $X=Spec(A),Y=Spec(B)$ with $S=Spec(C)$, with
						
						Then, $X\otimes_S Y$ exists in $Sch/S$, and is given by $Spec(A\otimes_C B)$ together with the ring morphism 
						\begin{gather*}
							A\rightarrow A\otimes B\qquad a\mapsto a\otimes 1\\
							B\rightarrow A\otimes B\qquad b\mapsto 1\otimes b
						\end{gather*}
					\end{lemma}
					\begin{proof}
						内容...
					\end{proof}
					
					\begin{corollary}
						For affine $X,Y,S$ and morphisms $f:X\rightarrow S,g:Y\rightarrow S$ then $X\times_S Y$ is affine.
					\end{corollary}
					\begin{proof}
						Omitted.
					\end{proof}
				\subsubsection{Fibre Products}
					\begin{Def}
						Given morphisms of schemes $f:X\rightarrow S$ and $g:Y\rightarrow S$, the \textbf{fibre product} is a scheme $X\times_S Y$ together with projection morphisms
						\begin{gather*}
							p:X\times_S Y\rightarrow X\qquad q:X\times_S Y\rightarrow Y
						\end{gather*}
						With the following commutative diagram.
						\begin{center}
							\begin{tikzcd}
								X\times_S Y\arrow[r,"q"]\arrow[d,"p"]&Y\arrow[d,"g"]\\
								X\arrow[r,"f"]&S
							\end{tikzcd}
						\end{center}
						which is universal among all diagrams of this sort.
					\end{Def}
					In other word, for $T\rightarrow X,T\rightarrow Y$, we have 
					\begin{center}
						\begin{tikzcd}
							T\arrow[rd,dashrightarrow]\arrow[rrd]\arrow[rdd]\\
							&X\times_S Y\arrow[r]\arrow[d]&Y\arrow[d]\\
							&X\arrow[r]&S
						\end{tikzcd}
					\end{center}
					\subsubsection{Open Subschemes}
					\begin{proposition}
						Let $f:X\rightarrow S$ and $g:Y\rightarrow S$ be morphisms of schemes with the same target, and $X\times_S Y,p,q$ be fibre product.
						
						 Suppose that $U\subseteq S,V\subseteq X,W\subseteq Y$ be open subschemes such that $f(V)\subseteq U,g(W)\subseteq U$.Then the canonical morphism
						 \begin{gather*}
						 	V\times_U W\rightarrow X\times_S Y
						 \end{gather*}
						 is an open immersion which identifies $V\times_U W$ with $p^{-1}(V)\cap q^{-1}(W)$.
					\end{proposition}
					\begin{proof}
						内容...
					\end{proof}
					
					\begin{proposition}
						There exists a correspondnece of points
						\begin{gather*}
							X\times_S Y\ni z\Leftrightarrow (x,y,s,p)
						\end{gather*}
						where:
						\begin{itemize}
							\item 
							\item 
						\end{itemize}
					\end{proposition}
					\begin{proof}
						\begin{center}
							\begin{tikzcd}
								X\times_S Y\\
								&Spec()&Spec&Y\\
								&Spec()&Spec\\
								&X&&S
							\end{tikzcd}
						\end{center}
					\end{proof}
				\subsubsection{Points in Fibre Product}
					\begin{proposition}
						内容...
					\end{proposition}
					\begin{proof}
						内容...
					\end{proof}
				\subsubsection{Immersions}
					\begin{lemma}
						内容...
					\end{lemma}
					\begin{proof}
						内容...
					\end{proof}
			\subsection{Base Change}
				\subsubsection{Base Change}
			
				\begin{Def}
					Let $S$ be a scheme:
					\begin{itemize}
						\item We say $X$ is a \textbf{scheme over $S$} to mean that $X$ is equipped with a morphisms of schemes 
						\begin{gather*}
							X\rightarrow S
						\end{gather*}
						The morphism $X\rightarrow S$ is sometimes called the \textbf{structure morphism};
						\item If $R$ is a ring we say $X$ is a \textbf{scheme over $S$} if $X$ is a scheme over $Spec(R)$;
						\item We denote $Mor_S(X,Y)$ the set of all morphisms from $X$ to $Y$ over $S$,
						\begin{center}
							\begin{tikzcd}
								X\arrow[rr]\arrow[rd]&&Y\arrow[ld]\\
								&S
							\end{tikzcd}
						\end{center}
						\item (Base Change of Scheme)
						\item (Base Change of Morphisms)
						\item 
						\item 
					\end{itemize}
				\end{Def}
				\subsubsection{Immersions}
				\begin{lemma}
					内容...
				\end{lemma}
				
				\begin{proof}
					内容...
				\end{proof}
				\subsubsection{Properties}
				\begin{Def}
					内容...
				\end{Def}
				\subsubsection{Scheme Theoretic Fibre}
				\begin{Def}
					内容...
				\end{Def}
				
				\begin{lemma}
					内容...
				\end{lemma}
				
				\begin{proof}
					内容...
				\end{proof}
			
			
			\subsection{Separation Axioms}
			
			'Hausdorffness for Schemes'
			
			\begin{Def}
				A topological space $X$ is Hausdorff iff $\Delta\subseteq X\times X$ is a closed subset.
				
				In anologue in AG is: Given a scheme $X$ over $S$
				\begin{gather*}
					\Delta_{X/S}:X\rightarrow X\times_S X
				\end{gather*}
				This is a unique morphism given by
				\begin{center}
					\begin{tikzcd}
						X\arrow[rd,"\Delta"]\arrow[rrrdd,"id_X",bend left=50]\arrow[dddrr,"id_X",bend right=50]\\
						&X\times X\arrow[rdd,"\pi_1",swap,bend right=20]\arrow[rrd,"\pi_2",bend left=20]\arrow[rd,dashrightarrow]\\
						&&X\times_S X\arrow[r,"pr_2"]\arrow[d,"pr_1"]&X\arrow[d]\\
						&&X\arrow[r]&S
					\end{tikzcd}
				\end{center}
				And $\Delta_{x/s}:X\stackrel{\Delta}{\rightarrow} X\times X\stackrel{p}{\rightarrow} X\times_S X$ is the unique morphism such that
				\begin{gather*}
					pr_1\circ \Delta_{X/S}=pr_2\circ \Delta_{X/S}
				\end{gather*}
			\end{Def}
			
			\begin{lemma}
				The diagonal morphisms of a morphism between affines is closed.
			\end{lemma}
			
			\begin{proof}
				内容...
			\end{proof}
			
			\begin{lemma}
				Let $X$ be a scheme over $S$. The diagonal morphsim $\Delta_{X/S}$ is an immersion.
			\end{lemma}
			
			\begin{proof}
				内容...
			\end{proof}
			
			\begin{Def}
				Let $f:X\rightarrow S$ be a morphism of schemes:
				
				\begin{itemize}
					\item We say $f$ is \textbf{separated} if the diagonal morphism $\Delta_{X/S}$ is a closed immersion;
					\item $f$ is \textbf{quasi-separated} if the diagonal morphism $\Delta_{X/S}$ is quasi-compact morphism;
					\item We say scheme $Y$ is \textbf{separated} if the morphism $Y\rightarrow Spec(Z)$ is separated;
					\item We say scheme $Y$ is \textbf{quasi-separted} if the morphism $Y\rightarrow Spec(Z)$ is quasi-separated;
				\end{itemize}
			\end{Def}
			
			\begin{lemma}
				内容...
			\end{lemma}
			
			\begin{proof}
				内容...
			\end{proof}
			
		\section{Quasi-Coherent Modules}
		\section{Schemes over Fields}
			\subsection{Schemes over a Field which is not Algebraically Closed}
				\subsubsection{title}
				\textbf{Construction} 
				
		\section{Local Properties of Schemes}
			\subsection{The Tangent Space}
				\subsubsection{Formal Derivatives}
		\section{Affine and Proper Morphisms}
			Affine Morphisms: \ref{Algebraic-Geometry-Affinization-Map}
	\chapter{Constructions}
		\section{Glueing Scheme,Relative Glueing, Relative Spectrum}
			\subsection{Glueing Scheme}
				\begin{Def}
					(Glueing Data) Let $I$ be a set. $(X_i,\mathcal{O}_i)$ be the locally ringed space. For each pair $i,j\in I$ and $U_{ij}\subseteq X_i$ be open subspace, and 
					\begin{gather*}
						\varphi_{ij}:U_{ij}\rightarrow U_{ji}
					\end{gather*}
				\end{Def}
				
				\begin{lemma}
					内容...
				\end{lemma}
				\begin{proof}
					内容...
				\end{proof}
			\subsection{Relative Spectrum via Glueing}
				Here: 
				
				\begin{itemize}
					\item $S$ is a scheme;
					\item $\mathcal{A}$ is a quasi-coherent $\mathcal{O}_S$-algebra;
				\end{itemize}
				\subsubsection{Relative Glueing}
					\begin{lemma}
						Let $S$ be a scheme, $\mathcal{B}$ be a basis for the topology of $S$. Suppose given the following data:
						\begin{itemize}
							\item For every $U\in \mathbb{B}$ a scheme
							\begin{gather*}
								f_U:X_U\rightarrow U
							\end{gather*}
							over $U$ of $S$;
							\item For $U,V\in \mathbb{B}$ a morphism $\rho^U_V:X_V\rightarrow X_U$ over $U$;
						\end{itemize}
						Assuming:
						\begin{itemize}
							\item Each $\rho_V^U$ induces an isomorphism
							\begin{gather*}
								X_V\rightarrow f_U^{-1}(V)
							\end{gather*}
							\item Whatever $W,V,U\in \mathcal{B}$, with $W\subseteq V\subseteq U$ have
							\begin{gather*}
								\rho^U_W=\rho^U_V\circ \rho^V_W
							\end{gather*}
						\end{itemize}
						Then, there exists a morphism $f:X\rightarrow S$ of schemes and isomorphism
						\begin{gather*}
							i_U:f^{-1}U\rightarrow X_U
						\end{gather*}
						over $U$ such that for $V,U\in \mathbf{B},V\subseteq U$ the composition
						\begin{gather*}
							X_V\stackrel{i_V^{-1}}{\rightarrow} f^{-1}(V)\stackrel{inclusion}{\rightarrow} f^{-1}(U)\stackrel{i_U}{\rightarrow} X_U
						\end{gather*}
						is the morphism $\rho^U_V$. Moreover $X$ is unique up to unique isomorphism over $S$.
					\end{lemma}
					\begin{proof}
						Firstly, we set
						\begin{gather*}
							F(T):=\{(g,\{h_U\}_{U\in \mathcal{B}})|g:T\rightarrow S,h_U:g^{-1}(U)\rightarrow X_U,f_U\circ h_U=g|g^{-1}(U),h_U|g^{-1}(V)=\rho^U_V\circ h_V,V\subseteq U\}
						\end{gather*}
						And we use the sheaf representability lemma:
						\begin{itemize}
							\item Sheaf Property:
							\item \begin{itemize}
								\item Representabillity
								\item Representable by Open Immersion:
								\item The collections $(F_i)$ covers $F$:
							\end{itemize}
						\end{itemize}
					\end{proof}
					
					\begin{lemma}
						Let $S$ be a scheme, $\mathcal{B}$ be a basis for the topology of $S$. Suppose given the following data:
						\begin{itemize}
							\item For every $U\in \mathcal{B}$ a scheme
							\begin{gather*}
								f_U:X_U\rightarrow U 
							\end{gather*}
							over $U$.
							\item For every $U\in \mathcal{B}$ a quasi-coherent sheaf $\mathcal{F}_U$ over $X_U$;
							\item For every $U,V\in \mathcal{B}$, morphism $\rho^U_W=\rho^U_V\circ \rho^V_W$;
							\item For every $U,V\in \mathcal{B}$, morphism $\theta^U_V:(\rho^U_V)^*\mathcal{F}_U\rightarrow \mathcal{F}_V$
						\end{itemize}
						assume that:
						\begin{itemize}
							\item 
							\item 
							\item 
							\item 
						\end{itemize}
						Then, there exists a morphism of schemes $f:X\rightarrow S$ together with a quasi-coherent sheaf on $X$ and isomorphism $i_U:f^{-1}(U)\rightarrow X_U;\theta_U:i^*_U\mathcal{F}_U\rightarrow \mathcal{F}|f^{-1}(U)$ and composition for $V\subseteq U$
						\begin{gather*}
							内容...
						\end{gather*}
					\end{lemma}
					\begin{proof}
						内容..
					\end{proof}
				\subsubsection{Affine Case}
				We will construct a scheme $X$ over $S$: Especially, if $S=Spec(R)$ is affine, and $A=\Gamma(S,\mathcal{A})$, then define:
				\begin{gather*}
					\underline{Spec}_S(\mathcal{A}):=Spec(A)
				\end{gather*} 
				With sheaf morphism induced by $R\rightarrow A$.
				\begin{gather*}
					\pi:Spec(A)=\underline{Spec}_S(\mathcal{A})\rightarrow S=Spec(R)
				\end{gather*}
				Furthermore: Suppose $\mathcal{A}$ be a quasi-coherent $\mathcal{O}_S$-alglebra, then we can construct the morphism of schemes
				\begin{gather*}
					\underline{Spec}_S(\mathcal{A})\rightarrow S
				\end{gather*}
				by lueing $Spec(\Gamma(U,\mathcal{A}))$ over the affine opens of $S$.
			
				\begin{lemma}
					Suppose $U\subseteq U'\subseteq S$ be affine opens, and
					\begin{gather*}
						A=\mathcal{A}(U),A'=\mathcal{A}(U')
					\end{gather*}
					The map of rings $A'\rightarrow A$ induces a morphism $Spec(A)\rightarrow Spec(A')$, and the diagram
					\begin{center}
						\begin{tikzcd}
							Spec(A)\arrow[r]\arrow[d]&Spec(A')\arrow[d]\\
							U\arrow[r]&U'
						\end{tikzcd}
					\end{center}
					is cartsian.
				\end{lemma}
				\begin{proof}
					内容...
				\end{proof}
				\begin{lemma}
					In situation above, suppose $U\subseteq U'\subseteq U''\subseteq S$ are affine opens, $A=\mathcal{A}(U),A'=\mathcal{A}(U'),A''=\mathcal{A}(U'')$, then the composition of morphisms $Spec(A)\rightarrow Spec(A'')$
				\end{lemma}
				\begin{proof}
					内容...
				\end{proof}
				\begin{lemma}
					There exists a morphism of schemes
					\begin{gather*}
						\pi:\underline{Spec}_S(\mathcal{A})\rightarrow S
					\end{gather*}
					with following properties:
					\begin{itemize}
						\item 
						\item 
					\end{itemize}
				\end{lemma}
				\begin{proof}
					内容...
				\end{proof}
			\subsection{Relative Spectrum as Functor}
				Short review: For $f:T\rightarrow S$, there exists a pullback for quasi-coherent sheaf of $\mathcal{O}_S$-algebras $\mathcal{A}$:
				\begin{gather*}
					f^*\mathcal{A}\in \mathcal{O}_T-Algs
				\end{gather*}
				Consider paisrs $(f:T\rightarrow S,\varphi)$ with corresponding morphism
				\begin{gather*}
					\varphi:f^*\mathcal{A}\rightarrow \mathcal{O}_T\\
					\varphi:f^{-1}\mathcal{A}\rightarrow \mathcal{O}_T\\
					\varphi:\mathcal{A}\rightarrow f_*\mathcal{O}_T		
				\end{gather*}
				Via $a:T'\rightarrow T$ we could get a second pair $(f'=f\circ a,\varphi'=a^*\varphi)$, which we call the \textbf{pullback}
				\begin{lemma}
					Let $F$ be the funtor associated to $(S,\mathcal{A})$ above, let $g:S'\rightarrow S$ be a morphism of schemes, $\mathcal{A}'=g^*\mathcal{A}$, let $F'$ be the functor associated to $(S',\mathcal{A}')$ above, there exists a canonnical isomorphism
					\begin{gather*}
						F'\cong h_{S'}\times_{h_S} F
					\end{gather*}
					of functors.
				\end{lemma}
				\begin{proof}
					内容...
				\end{proof}
				
				\begin{lemma}
					Let $F$ be the functor associated to $(S,\mathcal{A})$, if $S$ is affine, then $F$ is representable by the affine scheme $Spec(\Gamma(S,\mathcal{A}))$.
				\end{lemma}
				\begin{proof}
					内容...
				\end{proof}
				
				\begin{lemma}
					The functor 
				\end{lemma}
				\begin{proof}
					内容...
				\end{proof}
				\begin{lemma}
					内容...
				\end{lemma}
				\begin{proof}
					内容...
				\end{proof}
			\subsection{Relative Spectrum}
				\begin{Def}
					Consider $\pi:\underline{Spec}_S(\mathcal{A})\rightarrow S$, the \textbf{''universal family''} is a family of $\mathcal{O}_S$-algebras
					\begin{gather*}
						\mathcal{A}\rightarrow 
					\end{gather*}
				\end{Def}
		\section{Affine $n$-Space}
			\begin{Def}
				Let $S$ be a scheme and $n\geq0$. Define:
				\begin{gather*}
					A^n_S=(\mathcal{O}_S[T_1,\dots,T_n])
				\end{gather*}
				called \textbf{affine $n$-space over $S$}. 
				
				If $S=Spec(R)$ is affine, then this affine $n$-space over $R$ and we denote it $A^n_R$.
			\end{Def}
			
			Base Change: Especially, let $S\rightarrow S'$ of schemes we have
			\begin{gather*}
				g^*\mathcal{O}_S[T_1,\dots,T_n]=\mathcal{O}_{S'}[T_1,\dots,T_n]
			\end{gather*}
			
			and hence $A^n_{S'}=S'\times_S A_S^n$ is the base change. Therefore an alternative definition of affine $n$-space is the formula
			\begin{gather*}
				A^n_S=S\times_{Spec\ZZ} A^n_\ZZ
			\end{gather*}
			
			Also, a morphism from an $S$-scheme $f:X\rightarrow S$ to $A^n_S$ is given by a homomorphism of $\mathcal{O}_S$-algebras
			\begin{gather*}
				\mathcal{O}_S[T_1,\dots,T_n]\rightarrow f_*\mathcal{O}_X
			\end{gather*}
			
			with the bijection $Hom_{Sch/S}(X,\underline{Spec}_S(\mathcal{A}))\cong Hom_{\mathcal{O}_S-alg}(\mathcal{A},f_*\mathcal{O}_X)$We could define the value elementswise $h_i\in \Gamma(X,\mathcal{O}_X)$.
		\section{Vector Bundles}
			\begin{Def}
				Let $S$ be a scheme and $\epsilon$ be a quasi-coherent $\mathcal{O}_S$-module. The \textbf{vector bundle} associated to $\epsilon$ is
				\begin{gather*}
					V(\mathcal{E})=\underline{Spec}_S(Sym(\mathcal{E}))
				\end{gather*}
			\end{Def}
			Then, we have the grading
			\begin{gather*}
				\pi_*\mathcal{O}_{V(\mathcal{E})}=\oplus_{n\geq 0} Sym^n(\mathcal{E})
			\end{gather*}
			
			\begin{Def}
				Let $S$ be a scheme. A \textbf{vector bundle} $\pi:V\rightarrow S$ over $S$ is an affine morphism of schemes such that
			\end{Def}
			\subsection{Vector Bundle and Locally Free Bundles}
			\subsection{Torsors and Non-Abelian Cohomology	}
			\subsection{Vector Bundles and $GL_n$-Torsors}'
			\subsection{Flattening Stratification for Modules}
			\subsection{Divisors}
		\section{Cones}
		\section{Proj}
			\subsection{$Proj$ of a graded Ring(Algebra)}
			\subsection{Standard Open Covering}
					\begin{lemma}
						Let $S$ be a graded ring, and $f\in S$ wiht $deg(f)>0$:
						\begin{itemize}
							\item If $g\in S$ homogeneous of positive dgree with $D_+(g)\subseteq D_+(f)$, then:
								\begin{itemize}
									\item $f$ is invertible in $S_g$; $\frac{f^{deg(g)}}{g^{deg(f)}}$ is invertible in $S_{(g)}$;
									\item $g^e=af$ for some $e\geq 1$ and $a\in S$ homogeneous;
									\item there is a canonical $S$-algebra map $S_f\rightarrow S_g$;
									\item there is a canonical $S_0$-algebra map $S_{(f)}\rightarrow S_{(g)}$ compatible with the map $S_f\rightarrow S_g$;
									\item The map $S_{(f)}\rightarrow S_{(g)}$ induces an isomorphism
									\begin{gather*}
										(S_{(f)})_{g^{deg(f)}/f^{deg(g)}}\cong S_{(g)}
									\end{gather*}
									\item These maps induce a commutative diagram of topological spaces
									\begin{center}
										内容...
									\end{center}
									\item There are compatible canonical $S_f$ and $S_{(f)}$-module maps
									\begin{gather*}
										M_f\rightarrow M_g\qquad M_{(f)}\rightarrow M_{(g)}
									\end{gather*}
									
									induces an isomorphism
									\begin{gather*}
										(M_{(f)})_{g^{deg(f)}/f^{deg(g)}}=M_{(g)}
									\end{gather*}
								\end{itemize}
							\item Any open covering of $D_+(f)$ can be refined to be a finite open covering of the form
							\begin{gather*}
								D_+(f)=\cup_{i=1}^n D_+(g_i)
							\end{gather*}
							\item Let $g_1,\dots,g_n\in S$ be a homogeneous of positive degree. Then,
							\begin{gather*}
								D_+(f)\subseteq \cup D_+(g_i)\Leftrightarrow g_1^{deg(f)}/f^{deg(g_1)},\dots,g_n^{deg(f)}/f^{deg(g_n)}\mbox{ generate the unit ideal in $S(f)$}
							\end{gather*}
						\end{itemize}
					\end{lemma}
					\begin{proof}
						内容...
					\end{proof}
					
					\begin{Def}
						Let $S$ be a graded ring. Suppose that $D_+(f)\subseteq Proj(S)$ is a standard open. 
						
						A \textbf{standard open covering} of $D_+(f)$ is a covering $D_+(f)=\cup_{i=1}^n D_+(g_i)$ where $g_1,\dots,g_n\in S$ are homogeneous of positive degree.
					\end{Def}
				\subsection{Structure Sheaf}
			\subsection{Quas-Coherent Sheaves on $Proj$}
			\subsection{Invertible Sheaves on $Proj$}
			\subsection{Functoriality of $Proj$}
			\subsection{Morphisms into $Proj$}
		\section{Projective Space}
		\section{Relative Proj}
		\section{Projective Bundles}
		\section{Grassmannians}
	\chapter{Classification of Schemes}
		\section{Types of Schemes by Rings}
		\subsection{Local}
		\begin{Def}
			\begin{itemize}
				\item Let $P$ be a property of rings. We say that $P$ is \textbf{local} if the followings hold
				\begin{itemize}
					\item For any ring $R$, any $f\in R$ we have
					\begin{gather*}
						P(R)\Rightarrow P(R_f)
					\end{gather*}
					\item For any ring $R$ aand $f_i\in R,R=(f_1,\dots,f_n)$, then 
					\begin{gather*}
						\forall i,P(R_{f_i})\Rightarrow P(R)
					\end{gather*}
				\end{itemize}
				\item Let $P$ be a property of rings, $X$ be a scheme. We say $X$ is locally $P$, if for any $x\in X$ there exists an affine open neighbourhood $x\in U$ such that $\mathcal{O}_X(U)$ has the property $P$.
			\end{itemize}
		\end{Def}
		\begin{lemma}
			Let $X$ bea scheme, and $P$ is a local property of rings. The following are equivalent:
			\begin{itemize}
				\item 
				\item 
				\item 
				\item 
			\end{itemize}
		\end{lemma}
		\begin{proof}
			内容...
		\end{proof}
		
		\begin{lemma}
			Let $X$ be a scheme. Then $X$ is reduced iff $X$ is locally reduced.
		\end{lemma}
		\begin{proof}
			内容...
		\end{proof}
		
		\begin{lemma}
			The following properties of a ring $R$ are local:
			\begin{itemize}
				\item (Cohen-Macaulay)
				\item (Regular)
				\item (Absolutely Noetherian)
				\item 
			\end{itemize}
		\end{lemma}
		\begin{proof}
			内容...
		\end{proof}
		\section{Reduced Scheme}
			
			\begin{Def}
				Let $X$ be a scheme. We say $X$ is \textbf{reduced} if every local ring $\mathcal{O}_{X,x}$ is reduced.
			\end{Def}
			\subsubsection{Reduced on Sections and on Stalks}
			\begin{lemma}
				A scheme $X$ is reduced iff $\mathcal{O}_X(U)$ is a reduced ring for all open $U\subseteq X$ open.
			\end{lemma}
			
			\begin{proof}
				$\Rightarrow$ Let $f\in \mathcal{O}_X(U)$ such that $f^n=0$. Then, $f^n|u=0,\forall u\in U\Rightarrow f|u=0,\forall u\in U$. Thus $f=0$. 
				
				$\Leftarrow$ Conversely, if for all $\mathcal{O}_X(U)$, it is reduced. Then let $\bar{f}\in \mathcal{O}_{X,x}$ be nonzero with nonzero representive element $f\in \mathcal{O}(U)$. Thus $f^n\not=0\Rightarrow \bar{f}^n\not=0$
			\end{proof}
			\subsubsection{Affine Scheme}
			\begin{lemma}
				An affine scheme $Spec(R)$ is reduced iff $R$ is reduced.
			\end{lemma}
			\begin{proof}
				$\Rightarrow$ Take $\mathcal{O}_{Spec(R)}(Spec(R))$;
				
				$\Leftarrow$ If $R$ is reduced, then $\mathcal{O}_{Spec(R),p}=R_p$ is reduced.
			\end{proof}
			
			\subsection{Scheme Structure on Closed Subset}
			\begin{lemma}
				Let $X$ be a scheme, $T\subseteq X$ be a closed subset. There exists a unique closed subscheme $Z\subseteq X$ with the following properties:
				\begin{itemize}
					\item The underlying topological space of subscheme $|Z|=T$;
					\item $Z$ is reduced.
				\end{itemize}
			\end{lemma}
			
			\begin{proof}
				\begin{itemize}
					\item Firstly, consider the subscheme of $X$ by
					\begin{gather*}
						\mathcal{I}(U):=\{f\in \mathcal{O}_X(U)|\forall t\in T\cap U,f(t)=0\}
					\end{gather*}
					where $f(t):=\kappa(t)$. $\mathcal{I}$ is a sheaf on $X$:
					\begin{itemize}
						\item Identity: As a substructure of $X$;
						\item Pasting: Given $s_i\in \mathcal{I}(U_i)$. Then we will prove the pasting $s\in \mathcal{I}(U)$: Omitted.
					\end{itemize}
					Moreover, take the open affine $U=Spec(R)$ we have
					\begin{gather*}
						T\cap U=V(I)
					\end{gather*}
					for a \underline{unique} radical ideal $I\subseteq R$. Now given $p\in V(I)$ corresponding to $t\in T\cap U$ and an element . Hence:
					\begin{gather*}
						\mathcal{I}(U)=\cap_{p\in V(I)} p=I
					\end{gather*}
					Moreover, take standard open $D(g)\subseteq Spec(R)=U$ we have
					\begin{gather*}
						\mathcal{I}(D(g))=I_g
					\end{gather*}
					Thus $\mathcal{I}$ is quasi-coherent sheaf of ideals.
					\item Now take the closed subspace $Z$ associated to $\mathcal{I}$. For every affine open $U=Spec(R)$ of $X$ we see that
					\begin{gather*}
						Z\cap U=Spec(R/I)
					\end{gather*}
					where $I$ is a \underline{radical ideal}, so $Z$ is reduced. And the underlying topological space $|Z|=supp(\mathcal{O}_X/\mathcal{I})=T$.
					
					\item Uniqueness: Now let $Z'\subseteq X$ be a second closed subscheme with the properties. But notice that on every affines, we have $V(I)=V(I')$ because of uniqueness of $I$ or $I'$ we have $Z=Z'$.
				\end{itemize}
			\end{proof}
			
			\begin{Def}
				Let $X$ be a scheme, $Z\subseteq X$ be a closed subset. 
				
				A \textbf{scheme structure} on $Z$ is given by a closed subscheme $Z'$ of $X$ whose underlying set is equal to $Z$. 
				
				Espcially, the scheme structure on $X$ could be denoted as $X_{red}$.
			\end{Def}
			
			\textbf{Remark} For locally closed subset $T$, then $T$ is a closed subset of the open subscheme $X-\partial T$ of $X$, hence $\partial T=\bar{T}-T$ is the boundary of $T$ in the topological space $X$.
			
			($T$ is closed in $\bar{T}$, by definition)
			
			\begin{lemma}
				Let $X$ be a scheme. $Z\subseteq X$ be a closed subscheme. $Y$ be reuduced scheme. The morphism $f:Y\rightarrow X$ factors through $Z\Leftrightarrow f(Y)\subseteq Z$. 
				
				In particular, any morphism $Y\rightarrow X$ factors as
				\begin{gather*}
					Y\rightarrow X_{red}\rightarrow X
				\end{gather*}
			\end{lemma}
			\begin{proof}
				$\Rightarrow$ Omitted;
				
				$\Leftarrow$ Assume $f(Y)\subseteq Z$, then 
			\end{proof}
			
			\begin{lemma}
				The followings are equivalent
				\begin{itemize}
					\item 
					\item 
				\end{itemize}
			\end{lemma}
			\begin{proof}
				内容...
			\end{proof}
		\section{Integral, Irreducible and Reduced Schemes}
			\subsection{Integral Scheme}
			\subsection{Irreducible Scheme}
		
		\section{Noetherian Schemes}
				\begin{Def}
				\begin{itemize}
					\item We say $X$ is \textbf{locally Noetherian} if for every $x\in X$ has an affine open neighborhood $Spec(R)=U$ such that ring $R$ is Noetherian;
					\item We say $X$ is Noetherian if $X$ is locally Noetherian and quasi-compact;
				\end{itemize}
			\end{Def}
			\begin{lemma}
				(Equivalent Conditions for Noetherian Scheme) Let $X$ be a scheme, the following are equivalent:
				\begin{itemize}
					\item The scheme $X$ is locally Noetherian;
					\item For every affine open $U\subseteq X$ the ring $\mathcal{O}_X(U)$ is open;
					\item There exists an affine open covering $X=\cup U_i$ such that each each $\mathcal{O}_X(U_i)$ is Noetherian;
					\item There exists an affine covering $X=\cup X_j$ such that each $X_j$ is locally Noetherian.
				\end{itemize}
				Moreover, if $X$ is locally Noetherian then every open subschemme is locally Noetherian.
			\end{lemma}
			\begin{proof}
				内容...
			\end{proof}
			\subsubsection{Immersion}
			\begin{lemma}
				Any immersion $Z\rightarrow X$ with $X$ locally Noetherian is quasi-compact.
			\end{lemma}
			\begin{proof}
				内容...
			\end{proof}
			
			\begin{lemma}
				
			\end{lemma}
			\begin{proof}
				内容...
			\end{proof}
			
			\subsubsection{Underlying Space}
			\begin{lemma}
				A (locally) Noetherian scheme has a (locally) Noetherian underlying topological space.
			\end{lemma}
			\begin{proof}
				内容...
			\end{proof}
			\subsubsection{Morphism with Source Noetherian}
				\begin{lemma}
					内容...
				\end{lemma}
				\begin{proof}
					内容...
				\end{proof}
			\subsubsection{Closed Point}
			\subsubsection{Specialization}
				\begin{lemma}
					Let $X$ be a locally Noetherian scheme. Let $x'\rightsquigarrow x$ be a specialization of points of $X$. Then:
					\begin{itemize}
						\item There exists a discrete valuation ring $R$ and a morphism 
						\begin{gather*}
							f:Spec(R)\rightarrow X
						\end{gather*}
						such that the generic point $\eta$ of $Spec(R)$ maps to $x'$ and the special point maps to $x$; and
						\item provided $x\not=x'$ given a finitely generated field extension $K/\kappa(x')$, we may arrange it so that the extension $\kappa(\eta)/\kappa(x')$ induced by $f$ is isomorphic to the given one.
					\end{itemize}
				\end{lemma}
				\begin{proof}
					内容...
				\end{proof}
				
				\begin{lemma}
					Let $S$ be a Noetherian scheme. Let $T\subseteq S$ be an infinite subset. Then, there exists an infinite subset $T'\subseteq T$ such that there are no nontrivial specialization aong the points $T'$.
				\end{lemma}
				\begin{proof}
					内容...
				\end{proof}
				
				\begin{lemma}
					内容...
				\end{lemma}
				\begin{proof}
					内容...
				\end{proof}
				\subsubsection{Countable Dense Subset}
				\begin{lemma}
					Let $S$ be a Noetherian scheme, let $T\subseteq S$ be an infinite dense subset, then there exists a countable subset $E\subseteq T$ which is dense in $S$.
				\end{lemma}
				\begin{proof}
					内容...
				\end{proof}
		\section{Jacobson Schemes}
			\begin{Def}
				\begin{itemize}
					\item 	Let $X$ be a topological space, let $X_0$ be the set of closed points of $X$. We say: $X$ is \textbf{Jacobson} if every closed subset $Z\subseteq X$ is the closure of 
					\begin{gather*}
						Z\cap X_0
					\end{gather*}
					\item 	A scheme $S$ is said to be \textbf{Jacobson} if its underlying topological space is Jacobson.
				\end{itemize}
			\end{Def}
			\begin{lemma}
				An affine scheme $Spec(R)$ is Jacobson iff the ring $R$ is Jacobson.
			\end{lemma}
			\begin{proof}
				内容...
			\end{proof}
			\subsubsection{Equivalent Conditions}
			\begin{lemma}
				Let $X$ be a scheme. The following are equivalent
				\begin{itemize}
					\item The scheme $X$ is Jacobson;
					\item The scheme is 'locally Jacobson';(In the sense of 'local property' of Schemes by Rings)
					\item For every affine open $U\subseteq X$ the ring $\mathcal{O}_X(U)$ is Jacobson.
					\item There exists an affine open covering $X=\cup U_i$ such that each $\mathcal{O}_X(U_i)$ is Jacobson;
					\item There exists an open covering $X=\cup X_j$ such that each open subscheme $X_j$ is Jacobson.
				\end{itemize}
			\end{lemma}
			\begin{proof}
				内容...
			\end{proof}
			\subsubsection{Noetherian Jacobson Schemes}
			\begin{lemma}
				\begin{itemize}
					\item 
					\item 
					\item 
				\end{itemize}
			\end{lemma}
			\begin{proof}
				内容...
			\end{proof}
		\section{Normal Schemes}
			\begin{Def}
				A scheme $X$ is normal iff for all $x\in X$, the local ring $\mathcal{O}_{X,x}$ is a normal domain.
			\end{Def}
		\section{Cohen-Macaulay Schemes}
		\section{Regular Schemes}
		\section{Diemnsion of Schemes}
			\begin{Def}
				Let $X$ be a scheme:
				\begin{itemize}
					\item The dimension of $X$ is just the dimension of $X$ as a topological space;
					\item For $x\in X$ we denote 
					\begin{gather*}
						dim_x(X):=\min\{dim(U)|x\in U\mbox{ open}\}
					\end{gather*}
				\end{itemize}
			\end{Def}
			
			\begin{lemma}
				The followings are equal:
				\begin{itemize}
					\item The dimension of $X$;
					\item The superemum of the dimensions of local rings of $X$;
					\item The superemum of $dim_x(X)$ for $x\in X$;
				\end{itemize}
			\end{lemma}
			\begin{proof}
				\begin{itemize}
					\item 
					\item 
					\item 
				\end{itemize}
			\end{proof}
			
			\begin{lemma}
				Let $X$ be a scheme, $Y\subseteq X$ be an irreducible closed subset. Let $\xi\in Y$ be the generic point. Then
				\begin{gather*}
					codim(Y,X)=dim(\mathcal{O}_{X,\xi})
				\end{gather*}
				where the codimension is as defined in topology.
			\end{lemma}
			\begin{proof}
				内容...
			\end{proof}
			
			\begin{lemma}
				Let $X$ be a scheme with $x\in X$. Then $x$ is a generic point of an irreducible component of $X$ iff $dim(\mathcal{O}_{X,x})=0$
			\end{lemma}
			\begin{proof}
				内容...
			\end{proof}
			
			\subsection{Catenary Schemes}
				\begin{Def}
					\begin{itemize}
						\item A ring $A$ is called \textbf{catenary} if for any pair of prime ideals $p\subseteq q$ there exists a maximal chain of primes
						\begin{gather*}
							p=p_0\subseteq \dots\subseteq p_e=q
						\end{gather*}
						and all of these have the same length.
						\item Let $S$ be a scheme. We say $S$ is \textbf{catenary} if the underlying topological space of $S$ is catenary.
					\end{itemize}
				\end{Def}
				\begin{lemma}
					Let $S$ be a scheme, the following are equivalent:
					\begin{itemize}
						\item $S$ is catenary
						\item There exists an open covering of $S$ all of whose members are catenary;
						\item for every affine open $Spec(R)=U\subseteq S$ the ring $R$ is catenary;
						\item There exists an affine open covering $S=\cup U_i$ such that each $U_i$ is the spectrum of a catenary ring;
						\item For any $x\in X$, the local ring $\mathcal{O}_{X,x}$ is catenary.
					\end{itemize}
				\end{lemma}
				\begin{proof}
					内容...
				\end{proof}
				\subsubsection{Locally Noetherian Rings}
				\begin{lemma}
					Let $S$ be a locally Noetherian scheme. The following are equivalent:
					\begin{itemize}
						\item $S$ is catenary;
						\item locally in the Zariski topology there exists a dimension function on $S$;
					\end{itemize}
				\end{lemma}
				\begin{proof}
					内容...
				\end{proof}
			\subsection{Serre's Condition}
			\begin{Def}
				Let $X$ be a locally Noetherian scheme, $k\geq 0$:
				\begin{itemize}
					\item We say $X$ is regular in codimension $k$, or we say $X$ has the property $(R_k)$ if for every $x\in X$ we have
					\begin{gather*}
						dim(\mathcal{O}_{X,x})\leq k\Rightarrow \mathcal{O}_{X,x}\mbox{ is regular}
					\end{gather*}
					\item We say $X$ has property $(S_k)$ if for every $x\in X$ we have
					\begin{gather*}
						depth(\mathcal{O}_{X,x})\geq \min(k,dim(\mathcal{O}_{X,x}))
					\end{gather*}
				\end{itemize}
			\end{Def}
			\subsection{The Singular Locus}
		\section{Japanese and Nagata Schemes}
		\section{Quasi-Affine Scheme}
		\begin{Def}
			内容...
		\end{Def}
		\section{Ample Invertible Sheaves}
			\textbf{Review} Given an invertible shaef $\mathcal{L}$ on a locally ringed space $X$, and given a glocal section of $\mathcal{L}$, the set
			\begin{gather*}
				X_s:=\{x\in X|s\notin m_x \mathcal{L}_x\}
			\end{gather*}
			is open(Consider as open subscheme). Take the following remark
			\begin{gather*}
				X_{ss'}:=X_s\cap X_{s'}
			\end{gather*}
			where $ss'$ denote the section $s\otimes s'\in \Gamma(X,\mathcal{L}\otimes \mathcal{L}')$. 
			\begin{Def}
				Let $X$ be a scheme, $\mathcal{L}$ be an invertible $\mathcal{O}_X$-module. We say $\mathcal{L}$ is \textbf{ample} if
				\begin{itemize}
					\item $X$ is quasi-compact;
					\item for every $x\in X$ there exists an $n\geq 1$ and $s\in \Gamma(X,\mathcal{L}^{\otimes n})$ such that $x\in X_s$ and $X_s$ is affine.
				\end{itemize}
			\end{Def}
			\begin{lemma}
				Let $X$ be a scheme, $\mathcal{L}$ be an invertible $\mathcal{O}_X$-module, $n\geq 1$. 
				
				Then, $\mathcal{L}$ is ample iff $\mathcal{L}^{\otimes n}$ is ample.
			\end{lemma}
			\begin{proof}
				内容...
			\end{proof}
			
			\begin{lemma}
				Let $X$ be a scheme, $\mathcal{L}$ be an invertible $\mathcal{O}_X$-module. Let $s\in \Gamma(X,\mathcal{L})$, for any affine $U\cap X$ the intersection
				\begin{gather*}
					U\cap X_s
				\end{gather*}
				is affine.
			\end{lemma}
			\begin{proof}
				内容...
			\end{proof}
			\subsubsection{Tensor Product}
			\begin{lemma}
				Let $X$ bea scheme, let $\mathcal{L}$ and $\mathcal{M}$ be invertible $\mathcal{O}_X$-modules. If 
				\begin{itemize}
					\item 
					\item 
				\end{itemize}
			\end{lemma}
			\begin{proof}
				内容...
			\end{proof}
			\subsubsection{Topology}
		\subsection{Quasi-Coherent Sheaves and Ample Invertible Sheaves}
			\textbf{Situation} Let $X$ be a scheme. Let $\mathcal{L}$ be an ample invertible sheaf on $X$. Then, set
			\begin{gather*}
				S:=\Gamma_*(X,\mathcal{L})=\oplus_{n\geq 0} S_n,S_n:=\Gamma(X,\mathcal{L}^{\otimes n})
			\end{gather*}
			Set: $Y:=Proj(S)$, then $f:X\rightarrow Y$ be the canonical morphism
		\section{Appendix: How to find Suitable Affine Opens}
	\chapter{Morphsism of Schemes}
		\section{Classification of Properties}
			This section is just an elementary introduction.
			\subsection{Properties of Morphisms}
				\textbf{Notation} Let $f:X\rightarrow Y,g:Y'\rightarrow Y$. Now let 
				\begin{gather*}
					X_{Y'}:=X\times_Y Y'\qquad f_{Y'}:X_{Y'}\rightarrow Y'\qquad 
				\end{gather*}
				\begin{center}
					\begin{tikzcd}
						X_{Y'}\arrow[d,"g'"]\arrow[r,"f_{Y'}"]&Y'\arrow[d,"g"]\\
						X\arrow[r,"f"]&Y
					\end{tikzcd}
				\end{center}
				\begin{Def}
					Let $P$ be a property of morphisms of schemes 
					\begin{itemize}
						\item (COMP) It is stable under composition. That is: 
						\begin{center}
							For every (pair) $f:X\rightarrow Y,g:Y\rightarrow Z$ satisfies $P\Rightarrow g\circ f$ satisfies $P$;
						\end{center}
						\item (CANC) It is cancellable. That is :
						\begin{center}
							 For all $f:X\rightarrow Y,g:Y\rightarrow Z$; $g\circ f$ satisfies $P$$\Rightarrow$ $f$ satisfies $P$
						\end{center} 
						
						\item (BC) It is stable under base change, 
						\begin{center}
							For all $f:X\rightarrow Y,g:Y'\rightarrow Y$ satisfies $P$$\Rightarrow$  $f_Y':X_{Y'}\rightarrow Y'$.
						\end{center}
						\item (LOCT) It is (Zariski-)local on the target, if for all morphism $f:X\rightarrow Y$, all open covers $Y=\cup_i V_i$, the morphism $f$ satisfies $P$ iff $f_{V_i}:f^{-1}(V_i)\rightarrow V_i$ satisfies $P$ for all $i\in I$.
						\item (LOCS) It is (Zariski-)local on the source if for all morphism $f:X\rightarrow Y$ and all open covers $X=\cup U_i$, the moprhism $f$ satisfies $P$ iff $f|U_i$ satisfies $P$ for all $i\in I$.
						\item (Universally) A morphism $f:X\rightarrow Y$ is \textbf{universally} $P$ if for all morphism $g:Y'\rightarrow Y$, morphhism $f_{Y'}$ satisfies $P$.
					\end{itemize}
				\end{Def}
				
				\textbf{Remark: Universally} In algebra, we define
				\begin{Def}
					Ring homomorphism $\phi:A\rightarrow B$ is universally $P$: For any ring homomorphsim $A\rightarrow C$, the induced homomorphism $\phi_C:C\rightarrow B\otimes_A C$ satisfies $P$
					\begin{center}
						\begin{tikzcd}
							B\otimes_A C\arrow[r]\arrow[d]&C\arrow[d]\\
							B\arrow[r]&A
						\end{tikzcd}
					\end{center}
				\end{Def}
				
				\begin{example}
					Property 'flatness; is universal.
				\end{example}
				
				\begin{lemma}
					Let $f:X\rightarrow Y$ be aa morphism of schemes over $S$. Then: The following square is Cartesian:
					\begin{center}
						\begin{tikzcd}
							X\arrow[r,"\Gamma_f"]\arrow[d,"f"]&X\times_S Y\arrow[d,"f_Y"]\\
							Y\arrow[r,"\Delta_{Y/S}"]&Y\times_S Y
						\end{tikzcd}
					\end{center}
				\end{lemma}
				\begin{proof}
					内容...
				\end{proof}
				
				\begin{Def}
					If $P$ is a \underline{property of morphisms of schemes}. We say that a morphism satisfies $\Delta_P$ if $\Delta_{X/Y}$ satisfies $P$.
				\end{Def}
				
				\begin{proposition}
					内容...
				\end{proposition}
				\begin{proof}
					内容...
				\end{proof}
				
				\begin{lemma}
					内容...
				\end{lemma}
				\begin{proof}
					内容...
				\end{proof}
				
				\begin{theorem}
					(Cancellation Theorem)
				\end{theorem}
			\subsection{Topological Properties}
				\begin{Def}
					A morphism $(f,f^\#):(X,\mathcal{O}_X)\rightarrow (Y,\mathcal{O}_Y)$ of schemes:
					\begin{itemize}
						\item is inj./surj./bi. if $f$ is inj./surj./bi.;
						\item is open if $f$ is;
						\item is quasi-compact if $f$ is quasi-compact;
						\item is quasi-separated if $\Delta_{X/Y}:X\rightarrow X\times_Y X$ is quasi-compact;
						\item has finite fibers if $f^{-1}(y)$ is finite for all $y\in Y$;
					\end{itemize}
				\end{Def}
				
				\begin{proposition}
					The properties satisfies:
					\begin{itemize}
						\item 
						\item 
						\item 
						\item 
						\item 
						\item 
						\item 
						\item 
					\end{itemize}
				\end{proposition}
				\begin{proof}
						\begin{itemize}
						\item 
						\item 
						\item 
						\item 
						\item 
						\item 
						\item 
						\item 
					\end{itemize}
				\end{proof}
			\subsection{Scheme-Theoretic Properties}
				\begin{Def}
					\begin{itemize}
						\item (Locally Closed) Immersion
						\item Locally of finite type if: For all affine open 
						\item of finite type:
						\item quasi-finite if:
						\item affine: If
						\item Finite:
						\item Separated:
						\item Proper: If
					\end{itemize}
				\end{Def}
		\section{Valuative Criteria}
			\begin{Def}
				内容...
			\end{Def}
		\section{Quasi-Compact Morphism}
			\begin{Def}
				A morphism $f:X\rightarrow S$ is \textbf{quasi-compact} if its underlying map of topological space is quasi-compact. 
				
				i.e. for compact $K\subseteq S$ we have $f^{-1}(K)$ is compact.
			\end{Def}
			\begin{proposition}
				Let $f:X\rightarrow S$ be a morphism of schemes. The followings are equivalent:
				\begin{itemize}
					\item $f:X\rightarrow S$ is quasi-compact;
					\item the inverse image of every affine open is quasi-compact;
					\item there exists some affine open covering $S=\cup U_i$ such that $f^{-1}(U_i)$ is quasi-compact for all $i$;
				\end{itemize}
			\end{proposition}
			\begin{proof}
				内容...
			\end{proof}
			
			\begin{proposition}
				
			\end{proposition}
			\subsubsection{Closed Immersion}
			\begin{lemma}
				A closed immersion is quasi-compact.
			\end{lemma}
			\begin{proof}
				内容...
			\end{proof}
			\subsubsection{Specialization}
			
			\begin{proposition}
				Let $f:X\rightarrow S$ be a quasi-commpact morphism of schemes. The followings are equivalent
				\begin{itemize}
					\item $f(X)\subseteq S$ is closed;
					\item $f(X)\subseteq S$ is stable under specialization;
				\end{itemize}
			\end{proposition}
			\begin{proof}
				内容...
			\end{proof}
			\begin{proposition}
				Suppose $f:X\rightarrow S$ be a quasi-compact morphism of schemes. Then, $f$ is closed iff specialization lift along $f$.
			\end{proposition}
			\begin{proof}
				内容...
			\end{proof}
		\section{Universal Closedness}
			\begin{Def}
				A morphsim of schemes $f:X\rightarrow S$ is said to be \textbf{universally closed} if every base change
				\begin{gather*}
					f':X_{S'}\rightarrow S'
				\end{gather*}
				is closed.
			\end{Def}
			\subsection{Valuative Criterion}
		\section{Immersion}
			\subsection{Open Immersion}
			\subsection{Closed Immersion}
			\subsection{Closed Immersions and Quasi-Coherent Sheaaves}
			\subsection{Support}
		\section{Scheme Theoretic Image}
			\begin{lemma}
				Let $f:X\rightarrow Y$ be a morphism of schemes. There exists aclosed subscheme $Z\subseteq Y$ such that $f$ factors through $Z$, and that for any other closed subscheme $Z'\subseteq Y$ such that $f$ factors through $Z'$ we have $Z\subseteq Z'$.
			\end{lemma}
			\begin{proof}
				Define $\mathcal{I}:=Ker(\mathcal{O}_Y\rightarrow f_*\mathcal{O}_X)$, then by earlier discussion we have known, that $\mathcal{I}$ is quasi-coherent then we just take $Z$ to be the closed subscheme determined by $\mathcal{I}$.
				
				In general, to show that there exists a largest ideal $\mathcal{I}'$ contained in $\mathcal{I}$.
			\end{proof}
			
			\begin{Def}
				内容...
			\end{Def}
		\section{Scheme Theoretic Closure and Density}
		\section{Quasi-Compact Morphisms}
			Review: A map $f:X\rightarrow Y$ is called \textbf{quasi-compact} iff the preimage $f^{-1}(V)$ of every compact $V$ is compact.
			\begin{Def}
				A morphism of schemes if called \textbf{quasi-compact} if the underlying map of  topological maps is quasi-compact. 
			\end{Def}
			
			\begin{lemma}
				
			\end{lemma}
			\begin{proof}
				内容...
			\end{proof}
			\subsubsection{Quasi-Compact$\Rightarrow$ $f(X)$ is Closed$\leftrightarrow f(X)$   is stable under Spelization}
			\begin{lemma}
				内容...
			\end{lemma}
			\begin{proof}
				内容...
			\end{proof}
			\subsubsection{Quasi-Compact$\Rightarrow$ $f$ is closed$\leftrightarrow$ Specialization lift along $f$}
				\begin{lemma}
				内容...
			\end{lemma}
			
			\begin{proof}
				内容...
			\end{proof}
			
		\section{Monomorphisms}
		
			\begin{Def}
				A morphism of schemes is called a \textbf{monomorphism} if it is a monomorphism in the category of schemes.
			\end{Def}
			
			\begin{lemma}
				内容...
			\end{lemma}
			\begin{proof}
				内容...
			\end{proof}
		\section{Surjective Morphisms}
		\begin{Def}
			A morphism of schemes is said to be \textbf{surjective} if it is surjective on underlying topological spaces.
		\end{Def}
		\begin{lemma}
			The composition of surjective morphisms is surjective.
		\end{lemma}
		\begin{proof}
			Omitted.
		\end{proof}
		\subsection{Properties}
		\subsubsection{Fibre Product}
		\begin{lemma}
			Let $X,Y$ be schemes over a base schemes $S$. Then, given points $x\in X,y\in Y$, there is a point of $X\times_S Y$ mapping to $x$ and $y$ under the projection iff $x,y$ lie above the same point of $S$.
		\end{lemma}
		\begin{proof}
			内容...
		\end{proof}
		\subsubsection{The Base Change}
		\begin{lemma}
			内容...
		\end{lemma}
		\begin{proof}
			内容...
		\end{proof}
		\section{Radicial and Universally Injective Morphisms}
			\begin{Def}
				内容...
			\end{Def}
		
		\section{Dominant Morphisms}
			\begin{Def}
				A morphism $f:X\rightarrow S$ of schemes is called \textbf{dominant} if the image of $f$ is dense subset of $S$.
			\end{Def}
			\begin{example}
				Suppose $k$ be an infinite field, $(\lambda_i)_{i\in \NN}$ be the elements in $k$. Then, the morphism
				\begin{gather*}
					\sqcup_i Spec(k)\rightarrow Spec(k[x])
				\end{gather*}
				induced by $x\mapsto \prod_{i} (\lambda_i)$
			\end{example}
			\begin{lemma}
				Let $f:X\rightarrow S$ be a morphism of schemes, if every generic point of every irreducible component of $S$ is in the image of $f$, then $f$ is dominant.
			\end{lemma}
			\begin{proof}
				内容...
			\end{proof}
			
			\begin{lemma}
				Let $f:X\rightarrow S$ be \underline{quasi-coompact} morphism of schemes.
			\end{lemma}
			\begin{proof}
				内容...
			\end{proof}
			
			\subsection{Irreducible Component}
			\subsection{Base Change}
		\section{Affine Morphisms}
			\begin{Def}
				A morphism of scheme $f:X\rightarrow S$ is called \textbf{affine} if the inverse image of every affine open of $S$ is an affine open of $X$.
			\end{Def}
			\begin{lemma}
				Morphism $f:X\rightarrow S$ is affine $\Leftrightarrow f$ is separated and quasi-compact.
			\end{lemma}
			\begin{proof}
				内容...
			\end{proof}
			\subsection{Relative Spectrum}
				\begin{lemma}
					内容...
				\end{lemma}
				\begin{proof}
					内容...
				\end{proof}
				\begin{lemma}
					内容...
				\end{lemma}
				\begin{proof}
					内容...
				\end{proof}
		\section{Familes of Ample Invertible Modules}
			\begin{Def}
				Let $X$ be a scheme, $\{\mathcal{L}_i\}$ be a family of invertible $\mathcal{O}_X$-modules. We say $\{\mathcal{L}_i\}$ is an \textbf{ample family of invertible modules} on $X$ if:
				\begin{itemize}
					\item $X$ is quasi-compact;
					\item for every $x\in X$, there exists an $i\in I$ and $n\geq 1$, and $s\in \Gamma(X,\mathcal{L}_i^{\otimes n})$ such that
					\begin{itemize}
						\item $x\in X_s$;
						\item $X_s$ is affine;
					\end{itemize}
				\end{itemize}
			\end{Def}
		\section{Morphisms of Finite Type}
		\section{Dimension}
		\subsection{Morphisms and Dimension of Fibres}
			\begin{Def}
				 For $x\in X$ define
				 \begin{gather*}
				 	dim_x(X):=\min\{dim(U)|x\in U\mbox{ is open}\}
				 \end{gather*}
			\end{Def}
			
			\begin{lemma}
				Let $f:X\rightarrow S$ be a morphism of schemes, $x\in X$ and $s=f(x)$. Assume that $f$ is locally of finite type, then
				\begin{gather*}
					dim_x(X_S)=dim(\mathcal{O}_{X_s,x})+trdeg_{\kappa(s)}(\kappa(x))
				\end{gather*}
			\end{lemma}
			\begin{proof}
				内容...
			\end{proof}
		\subsection{Morphisms of Given Relative Dimension}
			\begin{Def}
				内容...
			\end{Def}
		\section{Ample}
			\subsection{Relatively Ample Sheaves}
			\subsection{Very Ample Sheaves}
			\subsection{Ample and Very Ample Sheaves Relative to Finite Type Morphisms}
		\section{Proper Morphisms}
		\section{Valuative Criteria}
		\section{Projective Morphisms}
		\section{Integral and Finite Morphisms}
		\section{Universal Homeomorphism}
			\subsection{Definitions}
			\subsection{Universal Homeomorphisms of Affine Schemes}
		\section{Normalization}
			\subsection{Absolute Weak Normalization and Seminormalization}
			\subsection{Relative Normalization}
			\subsection{Normalization}
			\subsection{Weak Normalization}
			\subsection{Weak Relative Noether Normalization}
		\section{Finite Locally Free Morphisms}
		\section{Rational Maps}
			\subsection{Definitions}
			\subsection{Birational Morphisms}
			\subsection{Zariski's Main Theorem}
		\section{Universally Bounded Fibres}
		\section{Sections of Quasi-Coherent Sheaves}
		
			\begin{lemma}
				Let $A$ be a ring. Let $I\subseteq A$ be a \underline{finitely generated ideal}. $M$ be an $A$-module.
				
				Then, there is a canonical map
				
				\begin{gather*}
					colim_{n}Hom_A(I^n,M)\rightarrow \Gamma(Spec(A)-V(I),\hat{M})
				\end{gather*}
				
				The maps is always injective.
				
				If for all $x\in M$, we have $Ix=0\Rightarrow x=0$  then this map is an isomorphism.
				
				In general, set $M_n:=\{x\in M|I^nx=0\}$, there exists an isomorphism
				
				\begin{gather*}
					colim_n Hom_A(I^n,M/M_n)\rightarrow \Gamma(Spec(A)-V(I),\hat{M})
				\end{gather*}
			\end{lemma}
			
			\begin{proof}
				内容...
			\end{proof}
			
			\begin{example}
				The first displayed map is not an isomorphism. 
				
				Let $k$ be a ring, 
			\end{example}
			
			\begin{lemma}
				Let $X$ be a quasi-compact scheme, $\mathcal{I}\subseteq \mathcal{O}_X$ be a quasi-coherent sheaf of ideals of finite type. Let $Z\subseteq X$ be the closed subscheme defined by $\mathcal{I}$ with $U=X-Z$. Let $\mathcal{F}$ be a quasi-coherent $\mathcal{O}_X$-module. The canonical map
				
				\begin{gather*}
					colim_n Hom_{\mathcal{O}_X} (\mathcal{I}^n,\mathcal{F})\rightarrow \Gamma(U,\mathcal{F})
				\end{gather*}
				
				is injective. Assume further that $X$ is quasi-separated, let $\mathcal{F}_n\subseteq \mathcal{F}$ be subsheaf of sections annihilated by $\mathcal{I}^n$, the canonical map
				
				\begin{gather*}
					colim_n Hom_{\mathcal{O}_X} (\mathcal{I}^n,\mathcal{F}/\mathcal{F}_n)\rightarrow \Gamma(U,\mathcal{F})
				\end{gather*}
			\end{lemma}
			
			\begin{proof}
				内容...
			\end{proof}
		\section{Thickenings}
			\subsection{Definitions}
			\subsection{Morphisms}
			\subsection{Picard Groups of Thinkenings}
			\subsection{Universal First Order Thickenings}
		\section{Unramified Morphisms}
			\subsection{Definitions}
			\subsection{Formally Unramified Morphisms}
		\section{Etale Morphisms}
			\subsection{title}
			\subsection{Formally Etale Morphisms}
		\section{Smooth Morphisms}
			\subsection{title}
			\subsection{Formally Smooth Morphisms}
			\subsection{Smoothness over a Noetherian Base}
		\section{The }
	\chapter{Classification of Schemes}
	\chapter{Divisors}
		\section{Associated Points}
			\begin{Def}
				Let $X$ be a scheme. Let $\mathcal{F}$ be a quasi-coherent sheaf on $X$:
				
				\begin{itemize}
					\item We say $x\in X$ is \textbf{associated to} $\mathcal{F}$ if the maximal ideal $m_x$ is associated to the $\mathcal{O}_{X,x}$-module $\mathcal{F}_x$, i.e. There exists $a\in \mathcal{F}_X$
					\begin{gather*}
						m_x=Ann(a)
					\end{gather*}
					\item Take the following notation:
					\begin{gather*}
						Ass(\mathcal{F})\qquad \mbox{or}\qquad Ass_X(\mathcal{F})
					\end{gather*}
					\item And: The associated points of $X$ are associated points of $\mathcal{O}_X$.
				\end{itemize}
			\end{Def}
			\textbf{Remark} These definitions are most useful when X is locally Noetherian and F of finite type.
			For example it may happen that a generic point of an irreducible component of X
			is not associated to X, see Example 2.7. In the non-Noetherian case it may be
			more convenient to use weakly associated points, see Section 5. Let us link the
			scheme theoretic notion with the algebraic notion on affine opens; note that this
			correspondence works perfectly only for locally Noetherian schemes.
			
			\begin{lemma}
				Let $X$ be a scheme. Let $\mathcal{F}$ be a quasi-coherent sheaf on $X$. Let $Spec(A)=U\subseteq X$ be an affine open, and set $M=\Gamma(U,\mathcal{F})$, $x\in U$, $p\subseteq A$ is the corresponding prime:
				\begin{itemize}
					\item If $p$ is associated to $M$, then $x$ is associated to $\mathcal{F}$;
					\item If $p$ is finitely generated, then the coverse holds as well.
				\end{itemize}
				In particular, if $X$ is Noetherian, then the equivalence
				\begin{gather*}
					p\in Ass(M)\Leftrightarrow x\in Ass(\mathcal{F})
				\end{gather*}
			\end{lemma}
			\begin{proof}
				内容...
			\end{proof}
			\subsection{Weakly Associated Points}
			\subsection{Morphisms and Weakly Associated Points}
			\subsection{Support}
			\begin{lemma}
				Let $X$ be a scheme, $\mathcal{F}$ be a quasi-coherent $\mathcal{O}_X$-module, then
				\begin{gather*}
					Ass(\mathcal{F})\subseteq Supp(\mathcal{F})
				\end{gather*}
			\end{lemma}
			\begin{proof}
				内容...
			\end{proof}
			\subsubsection{Short Exact Sequence}
			\subsubsection{Locally Noetherian Scheme}
			\subsection{Morphisms and Associated Points}
			\subsection{Embedded Points}
		\section{Assassin}
			\subsection{Relative Assassin}
			\subsection{Relative Weak Assassin}
		\section{Fitting Ideals}
		\section{Effective Cartier Divisors}
			\subsection{Definitions}
			\subsection{Effective Cartier Divisors and }
		\section{Norm}
		\section{Weil Divisors}
	\chapter{Varities}
		\section{Definitions}
			\begin{Def}
				Let $k$ be a field. A \textbf{variety} is a scheme $X$ over $k$ such that $X$ is integral and the structure morphism
				\begin{gather*}
					X\rightarrow Spec(k)
				\end{gather*}
				
				is separated and of finite type.
			\end{Def}
			\subsection{Varieties and Rational Maps}
		\section{Base Field Changes}
	\chapter{Topology of Schemes}
	\chapter{Limits of Schemes}
	\chapter{Chow Homology and Chern Classes}
	\chapter{Intersection Theory}
		\textbf{Remark} 
		\section{Cycles}
			\subsection{Definitions}
	\chapter{title}
\end{document}